\section{Shared results VII: unramified depth, Bessel models and moduli}\label{sec:pf:sh-local}

\providecommand{\shlocalpsi}{\psi_{\mathrm{std}}}
\providecommand{\shlocalHS}[1]{\|#1\|_{\mathrm{HS}}}
\providecommand{\shlocalE}{\mathcal{E}}
\providecommand{\shlocalOE}{\OO_{E}}
\providecommand{\shlocalZOE}{\Z_p\otimes\OO_E}

This section proves twelve results of the formal development \cite{LeanSource2026} that are cited by
earlier sections of Part~II. They fall into four groups.
\begin{enumerate}[label=(\arabic*)]
\item \emph{The real place.} The uniqueness of Bessel maps for discrete series whose Harish-Chandra parameter
has a large gap (\cref{pf:sh-local:bmu,pf:sh-local:uniqgap}), used in
\cref{sec:pf:open0,sec:pf:stageU,sec:pf:thmS-carrier,sec:pf:thmS,sec:pf:fs}; the energy bound for the real
objects at infinity of a Hida setting (\cref{pf:sh-local:realobj}), used in
\cref{sec:pf:open0,sec:pf:stageU,sec:pf:thmS-leaves}; and Harish-Chandra's finiteness theorem in the form
needed for the row spaces (\cref{pf:sh-local:findim}), used in \cref{sec:pf:stageU}.
\item \emph{Isomorphisms $\C\cong\C_q$} extending a given embedding of $\Qbar$ (\cref{pf:sh-local:ringequiv}),
used in \cref{sec:pf:reckernel,sec:pf:rows}.
\item \emph{Unramified depth.} The places of the host field above the small primes
(\cref{pf:sh-local:twoplaces}), used in \cref{sec:pf:lift1,sec:pf:onto2}; three statements on characters of
formal groups with jet charts, the Hodge--Tate inclusion $(\mathrm H_p)$, the unramified bound $(\mathrm U_p)$
and the split-place bound $(\mathrm A_p)$ (\cref{pf:sh-local:Hp,pf:sh-local:Up,pf:sh-local:Ap}), used in
\cref{sec:pf:onto1,sec:pf:onto2}; and the period tower theorem (\cref{pf:sh-local:ptd}), used in
\cref{sec:pf:tower}.
\item \emph{Moduli.} The existence of unitary moduli schemes at every level (\cref{pf:sh-local:moduli}) and the
theta-constant formula (\cref{pf:sh-local:key}), used in \cref{sec:pf:deepstr}.
\end{enumerate}
Several of these results are formalised versions of classical theorems: the existence of field isomorphisms
$\C\cong\C_q$ (Steinitz), the archimedean multiplicity-one theorem for Bessel models
\cite{AizenbudGourevitchRallisSchiffmann2010,SunZhu2012,JiangSunZhu2010}, Harish-Chandra's finiteness theorem
\cite{HarishChandra1968}, the existence of fine moduli schemes of polarised abelian varieties with level structure
\cite{MumfordFogartyKirwan1994} and Moret-Bailly's \emph{formule cl\'e} \cite{MoretBailly1985}. For each of them we
state precisely what is proved and how it compares with the published statement; the formal proofs are
independent of the literature, and we give them at the level of their main steps. The arguments that carry the
analytic weight at the real place (the orbit multiplicity argument for Bessel maps and the abstract energy bound,
\cref{pf:sh-local:orbitmult,pf:sh-local:energyabs}) and the algebraic number theory of the places of the host
(\cref{sec:pf:sh-local:host}) are given in full. Deeper lemmas of the formal proofs (decompositions of unitary
groups, length bounds for lattices over valuation rings, constructions with biextensions and theta functions) are
cited by their labels, each with a statement of what is used.

Throughout, $\vol$, Haar measures and integrals are those of the formal development; ``Hilbert space'' means a
complex inner product space, complete unless said otherwise; inner products $\langle x,y\rangle$ are conjugate-linear in
$x$ and linear in $y$ (the convention of the formal development); for vectors $v,w$ of inner product spaces
$v\otimes w^*$ denotes the rank-one operator $x\mapsto\langle w,x\rangle v$; and $\shlocalHS{A}^2=\sum_i\|Ae_i\|^2$ for an
orthonormal basis $(e_i)$ of the (finite-dimensional) source of $A$.

% ---------------------------------------------------------------------------
\subsection{Uniqueness of Bessel maps in large gap}\label{sec:pf:sh-local:bessel}

\subsubsection{The setting}
Fix integers $k\ge1$, $r\ge0$ and put $n=2r+k+1$. Index the coordinates of $\C^n$ by $\{0,\dots,n-1\}$ and let
$(e_i)$ be the standard basis. Let $H=H_{k,r}$ be the hermitian matrix with entries $H_{a,n-1-a}=H_{n-1-a,a}=1$
for $a<r$, $H_{ii}=1$ for $r\le i<r+k$, $H_{r+k,r+k}=-1$, and all other entries $0$; write
$\langle x,y\rangle_H=y^*Hx$. For $a<r$ put $e_a^*=e_{n-1-a}$; the vectors $e_a,e_a^*$ span $r$ hyperbolic planes,
and their orthogonal complement is $W_0=\langle e_r,\dots,e_{r+k}\rangle$, on which the form is
$\mathrm{diag}(1,\dots,1,-1)$. Let $G=\Uni(H)=\{g\in M_n(\C):g^*Hg=H\}$, with Lie algebra
$\mathfrak g=\{X:X^*H+HX=0\}$.

Let $z\in W_0$ with $T(z):=-\langle z,z\rangle_H>0$. Put
\[
U_z=\{g\in G:\ ge_a=e_a,\ ge_a^*=e_a^*\ (a<r),\ gz=z\},
\]
\[
N=\{g\in G:\ g_{ia}=\delta_{ia}\ \text{for }a<r,\ i\ge a;\quad g_{ij}=\delta_{ij}\ \text{for }r\le j\le r+k,\ i\ge r\}.
\]
Thus $N$ consists of the elements of $G$ which are unipotent upper triangular on the isotropic flag
$\langle e_0\rangle\subset\langle e_0,e_1\rangle\subset\dots\subset\langle e_0,\dots,e_{r-1}\rangle$ and act
trivially on $W_0$ modulo this flag. For $c\in\C^r$ and $g\in M_n(\C)$ let
\[
\ell_c(g)=\sum_{a<r-1}c_a\,g_{a,a+1}+c_{r-1}(gz)_{r-1},\qquad \psi_c(g)=\exp\bigl(i\Re\ell_c(g)\bigr),
\]
($\ell_c=0$ if $r=0$), and let $\shlocalpsi(g)=\exp\bigl(i\Re(-2i\,\ell_{\mathbf 1}(g))\bigr)$ be the
\emph{standard character}, $\mathbf 1=(1,\dots,1)$. A direct computation shows that $\ell_c$ is additive on $N$, so that
$\psi_c$ and $\shlocalpsi$ restrict to characters of $N$ (the generic characters of the Bessel subgroup).

A \emph{representation} $\sigma$ of a closed matrix group $\Gamma\subseteq M_n(\C)$ on a Hilbert space
$\mathcal H$ is a map $g\mapsto\sigma(g)$ into the bounded operators with $\sigma(1)=1$ and
$\sigma(gg')=\sigma(g)\sigma(g')$ on $\Gamma$; it is \emph{unitary} if every $\sigma(g)$, $g\in\Gamma$, is
isometric and $g\mapsto\sigma(g)v$ is continuous on $\Gamma$ for every $v$. A vector $v$ is \emph{smooth} if
$g\mapsto\sigma(g)v$ coincides near every point of $\Gamma$ with a $C^\infty$ map defined on an open set of
$M_n(\C)$; we write $\mathcal H^\infty$ for the space of smooth vectors. For $X$ in the Lie algebra and smooth
$v$, $d\sigma(X)v=\frac{d}{dt}\sigma(\exp tX)v|_{t=0}$, and for a word $W=(X_1,\dots,X_l)$ we write
$d\sigma(W)=d\sigma(X_1)\cdots d\sigma(X_l)$.

\begin{definition}[Bessel maps]\label{pf:sh-local:def:bessel}
Let $\sigma$ be a unitary representation of $G$ on $\mathcal H$, $\pi$ a representation of $U_z$ on a Hilbert
space $V$ and $\psi$ a function on $G$. A \emph{smooth Bessel map} for $(z,\psi,\sigma,\pi)$ is a map
$B:\mathcal H\to V$ which is additive and $\C$-linear on $\mathcal H^\infty$, satisfies
\[
B(\sigma(hm)v)=\psi(m)\,\pi(h)B(v)\qquad(h\in U_z,\ m\in N,\ v\in\mathcal H^\infty),
\]
and is \emph{Sobolev-continuous}: there are $C$ and a finite list $\mathcal W$ of words in $\mathfrak g$ with
$\|Bv\|\le C\sum_{W\in\mathcal W}\|d\sigma(W)v\|$ for all $v\in\mathcal H^\infty$. A map with the first and the
last property (but not necessarily the equivariance) is called \emph{Sobolev-linear}.
\formal{PI.BesselMapUniq.IsBesselMapSm}
\end{definition}

\begin{definition}[Discrete series, gap, Harish-Chandra type]\label{pf:sh-local:def:ds}
For $\eps\in\{\pm1\}^n$ let $J_\eps=\mathrm{diag}(\eps)$. A vector $\lambda\in\R^n$ is a \emph{Harish-Chandra
parameter} if it is strictly decreasing and $\lambda_a-\frac{n-1}2\in\Z$ for all $a$; the pair $(\eps,\lambda)$
has \emph{gap at least $h$} if $|\lambda_a-\lambda_b|\ge h$ whenever $\eps_a\ne\eps_b$. A representation $\sigma$
of $G$ is a \emph{discrete series with Harish-Chandra parameter $(\eps,\lambda)$} if for some unitary matrix $P$
with $P^*J_\eps P=H$ the representation $g\mapsto\sigma(P^{-1}gP)$ of $\Uni(J_\eps)$ is a discrete series with
parameter $\lambda$ in the following sense: its space is complete and non-zero, it is unitary, topologically
irreducible and square integrable (all matrix coefficients lie in $L^2$ of a Haar measure), it contains the
$K$-type ($K=\Uni(J_\eps)\cap\Uni(n)$) of highest weight $\lambda+\rho_n-\rho_c$, and every $K$-type $\mu$ which
occurs satisfies $|\lambda+\rho_n+\rho_c|\le|\mu+2\rho_c|$. A representation $\pi$ of $U_z$ on $V$ has
\emph{Harish-Chandra type $u\in\R^k$} if for some $g\in G$ fixing all $e_a,e_a^*$ with
$g(\sqrt{T(z)}\,e_{r+k})=z$ the representation $h\mapsto\pi(g\iota(h)g^{-1})$ of $\Uni(k)$ is non-zero,
finite-dimensional, unitary and consists of the single $K$-type with parameter $u+\rho_n-\rho_c$ (for the
compact group $\Uni(k)$); here $\iota(h)$ acts by $h$ on $\langle e_r,\dots,e_{r+k-1}\rangle$ and trivially on
the other basis vectors.
\formal{PI.UpqRep.IsDSHpk}\formal{PI.UpqRep.IsHCTypeW}
\end{definition}

\subsubsection{The theorem}

\begin{theorem}[Uniqueness of Bessel maps in large gap]\label{pf:sh-local:bmu}
Let $k\ge1$, $r\ge0$, $n=2r+k+1$, let $(\eps,\lambda)$ be a pair with $\lambda$ a Harish-Chandra parameter and
gap at least $h\ge120\,n^3$, and $\sigma$ a discrete series of $G=\Uni(H_{k,r})$ on $\mathcal H$ with
Harish-Chandra parameter $(\eps,\lambda)$. Let $z\in W_0$ with $T(z)>0$, $c\in(\C^\times)^r$, and $\pi$ a
representation of $U_z$ on a Hilbert space $V$ of Harish-Chandra type $u$ for some $u\in\R^k$. Then any two
smooth Bessel maps $B_1,B_2$ for $(z,\psi_c,\sigma,\pi)$ are linearly dependent on smooth vectors: there are
$(a,b)\in\C^2\setminus\{0\}$ with $aB_1(v)+bB_2(v)=0$ for all $v\in\mathcal H^\infty$.
\formallabel{CompanionTheory.Bmu.besselMapUniqGap}
\formallabel{CompanionTheory.Bmu.uniq_of}
\end{theorem}

\begin{corollary}\label{pf:sh-local:uniqgap}
The proposition $\sfname{UniqGap}$ holds; it is literally the statement of \cref{pf:sh-local:bmu}.
\formallabel{CompanionTheory.Lime.BM.uniqGap_holds}
\end{corollary}

\begin{proof}
The two propositions $\sfname{UniqGap}$ and $\sfname{BesselMapUniqGap}$ have identical statements (they are
formulated with two names for the same type of $n\times n$ matrices), so \cref{pf:sh-local:bmu} applies.
\end{proof}

\begin{remark}[Relation to the literature]\label{pf:sh-local:rem:bmu}
A smooth Bessel map $B$ defines the form $(v,w)\mapsto\langle w,B(v)\rangle$ on
$\mathcal H^\infty\times V$, linear in $v$, which is continuous for the Fr\'echet topology of $\mathcal H^\infty$ (by
Sobolev-continuity) and transforms under $U_z\ltimes N$ by $\psi_c\otimes\pi$; thus $B$ gives an element of
$\Hom_{U_z\ltimes N}(\mathcal H^\infty\otimes\overline{V},\psi_c)$, where $\mathcal H^\infty$ is the
Casselman--Wallach globalisation of the Harish-Chandra module of $\sigma$ \cite{Wallach1992}. The pair
$(G,U_z\ltimes N)$ with the generic character $\psi_c$ is a Bessel pair of unitary groups in the sense of
Gan--Gross--Prasad, and the uniqueness of Bessel models \cite[Theorem~A]{JiangSunZhu2010} (which rests on the
multiplicity-one theorem \cite{AizenbudGourevitchRallisSchiffmann2010,SunZhu2012}) asserts that this space has
dimension at most one for all irreducible Casselman--Wallach representations of $G$ and $U_z$. Hence
\cref{pf:sh-local:bmu} is a special case of the published theorem: it is restricted to discrete series with gap at
least $120n^3$ and to Bessel maps with values in a finite-dimensional unitary representation of the compact group
$U_z\cong\Uni(k)$. The formal proof does not use these references; it is given below. The gap hypothesis enters
only through the decay of matrix coefficients (\lref{CompanionTheory.FamilyDerivativeExpansion.DecS.decayM}).
\end{remark}

\subsubsection{Inputs proved elsewhere}
The proof uses the following results of the development.
\begin{itemize}
\item \emph{Decay} (\lref{CompanionTheory.FamilyDerivativeExpansion.DecS.decayM}): if $\sigma$ is a discrete
series with Harish-Chandra parameter $(\eps,\lambda)$ of gap $h$ and $\theta^*(\eps)+m<h$, where $\theta^*(\eps)$ is
an explicit threshold, then for every smooth $v_0$ there are $C$ and finitely many words $W$ such that
$|\langle v_0,\sigma(x)w\rangle|\le C\|x\|^{-b}\sum_W\|d\sigma(W)w\|$ for all $b\le m$, all smooth $w$ and all
$x\in G$, where $\|x\|$ is the Frobenius norm. Moreover $\theta^*(\eps)\le40n^3$
(\lref{CompanionTheory.FamilyDerivativeExpansion.DecS.thr_le}).
\item \emph{Basic properties of discrete series}
(\lref{GlobalTheory.T.discreteSeries_unitary_smooth_basic}, with
\lref{CompanionTheory.D.gk.admitted.k_finite_smooth}): $G$-translates of smooth vectors are smooth,
$K$-finite vectors are smooth, and there is a $K$-finite unit vector $w_0$ in the lowest $K$-type.
\item \emph{Square integrability and irreducibility}
(\lref{ReciprocityTheory.ChartRestriction.DS.Mult1.TopTA.sqInt_of_DSHpk},
\lref{ReciprocityTheory.ChartRestriction.DS.Mult1.TopTA.irr_of_DSHpk}): $\sigma$ is unitary and irreducible on
$G$ and its matrix coefficients are square integrable for every Haar measure; and the properties of $\pi$
(\lref{ReciprocityTheory.ChartRestriction.DS.Mult1.TopTA.pi_facts}): $\pi$ is a continuous unitary irreducible
representation of $U_z$.
\item \emph{Diagonal data and Schur orthogonality}
(\lref{ReciprocityTheory.ChartRestriction.DS.Mult1.stDiagExists},
\lref{ReciprocityTheory.ChartRestriction.DS.Mult1.stSchurG}): for $z,\pi$ as in the theorem there are Haar
measures $\mu_G$ on $G$, $\mu_N$ on $N$ and a Haar probability measure $\mu_W$ on $U_z$ (a \emph{diagonal real
place datum} $\mathcal D$ for $(z,\pi)$); and for every irreducible unitary representation of $G$ with square
integrable coefficients there is $d>0$ (the formal degree) with
\begin{equation}\label{eq:pf:sh-local:schur}
\int_G\langle b,\sigma(x)a\rangle\overline{\langle b',\sigma(x)a'\rangle}\,d\mu_G(x)=\frac{\langle a',a\rangle\langle b,b'\rangle}d .
\end{equation}
\item \emph{The multiplicity-one map} (\lref{ReciprocityTheory.ChartRestriction.DS.Mult1.thmM}): let
$\mathcal D$ be a diagonal real place datum for $(z,\pi)$, and $X$ an irreducible unitary representation $\sigma$ of
$G$ on $\mathcal H$ with square integrable coefficients and a unit vector $w$ such that
$|\langle w,\sigma(x)w\rangle|\le C_w\|x\|^{-a}$ with $a\ge4n^3$. Let $\mathcal V$ be the space of $v$ with
$|\langle w,\sigma(x)v\rangle|\le C\|x\|^{-a}$ for some $C$. Then there is a map $l:\mathcal H\to V$, additive and
linear on $\mathcal V$, with $l(\sigma(mh)v)=\shlocalpsi(m)\pi(h)l(v)$ for $h\in U_z$, $m\in N$, $v\in\mathcal V$,
such that for $v,v'\in\mathcal V$ the \emph{Bessel operator}
\begin{equation}\label{eq:pf:sh-local:beop}
\mathfrak B(v,v')=\int_{U_z\times N}\langle v',\sigma(mh)v\rangle\,\overline{\shlocalpsi(m)}\,\pi(h)^{-1}\,d\mu_W(h)\,d\mu_N(m)
\end{equation}
equals the rank-one operator $l(v)\otimes l(v')^*$.
\item \emph{Cyclicity} (\lref{CompanionTheory.D.Upq.Sm.cyc_of_smooth}): if $w\ne0$ is smooth, then for every
smooth $x$, every finite list of words $\mathcal W$ and every $\delta>0$ there are $c_i\in\C$ and $g_i\in G$ with
$\sum_{W\in\mathcal W}\|d\sigma(W)(x-\sum_ic_i\sigma(g_i)w)\|<\delta$.
\end{itemize}

\subsubsection{Lemmas}

\begin{lemma}[Torus conjugacy]\label{pf:sh-local:torconj}
Let $z\in W_0$ and $c\in(\C^\times)^r$. There is $\tau\in G$ commuting with every element of $U_z$ such that for
every $m\in N$ one has $\tau m\tau^{-1}\in N$ and $\psi_c(\tau m\tau^{-1})=\shlocalpsi(m)$.
\formallabel{GlobalTheory.T.TorusConjugacy}
\end{lemma}

\begin{proof}
Define $T_{r-1}=-2i/c_{r-1}$ and recursively $T_a=-2i\,T_{a+1}/c_a$ for $a<r-1$; all $T_a$ are non-zero. Let
$\tau$ be the diagonal matrix with entries $T_a$ at $a<r$, $1$ at $r\le i\le r+k$ and $\overline{T_a}^{-1}$ at
$n-1-a$. Then $\overline{\tau_{ii}}\tau_{jj}=1$ whenever $H_{ij}\ne0$, so $\tau\in G$. An element $h\in U_z$ fixes
$e_a,e_a^*$ and preserves $H$, hence preserves their orthogonal complement $W_0$; so $h$ is the identity outside
$W_0$ while $\tau$ is the identity on $W_0$, and $\tau h=h\tau$. Conjugation by the diagonal matrix $\tau$
multiplies the entry $(i,j)$ by $\tau_{ii}/\tau_{jj}$; it preserves the conditions defining $N$ (they prescribe
diagonal entries and the vanishing of entries), so $\tau N\tau^{-1}=N$. Finally
$(\tau m\tau^{-1})_{a,a+1}=(T_a/T_{a+1})m_{a,a+1}$ and, as $\tau^{-1}z=z$,
$(\tau m\tau^{-1}z)_{r-1}=T_{r-1}(mz)_{r-1}$; by the choice of the $T_a$ this gives
$\ell_c(\tau m\tau^{-1})=-2i\,\ell_{\mathbf 1}(m)$, hence $\psi_c(\tau m\tau^{-1})=\shlocalpsi(m)$.
\end{proof}

The next three lemmas hold for an arbitrary closed matrix group $\Uni(H_m)=\{g:g^*H_mg=H_m\}$ with $H_m^2=1$ and a
unitary representation $\sigma$ on a Hilbert space $\mathcal H$; $V$ denotes a finite-dimensional Hilbert space.

\begin{lemma}[Factorisation]\label{pf:sh-local:factor}
If $B:\mathcal H\to V$ is Sobolev-linear, there are finitely many words $W_1,\dots,W_N$ in the Lie algebra and
bounded operators $\beta_i:\mathcal H\to V$ with $B(v)=\sum_i\beta_i(d\sigma(W_i)v)$ for all smooth $v$.
\formallabel{GlobalTheory.T.factorization'}
\end{lemma}

\begin{proof}
Let $\mathcal W=(W_1,\dots,W_N)$ be the words of the Sobolev bound $\|Bv\|\le C\sum_i\|d\sigma(W_i)v\|$ and
$J:\mathcal H^\infty\to\mathcal H^N$, $Jv=(d\sigma(W_i)v)_i$, a linear map. Then $\|Bv\|\le|C|N\|Jv\|$, so
$B=T_0\circ J$ for a bounded linear map $T_0$ on the subspace $J(\mathcal H^\infty)$, which extends (as $V$ is
finite-dimensional) to a bounded operator $T:\mathcal H^N\to V$
(\lref{GlobalTheory.CompanionTheory.BmuFactor.exists_factor}). Put
$\beta_i=T\circ\iota_i$ with $\iota_i$ the $i$-th coordinate embedding.
\end{proof}

\begin{lemma}[Translation of derivatives]\label{pf:sh-local:dsad}
For $y\in\Uni(H_m)$, a word $W=(X_1,\dots,X_l)$ in the Lie algebra and a smooth vector $u$, the word
$y^{-1}Wy=(y^{-1}X_1y,\dots,y^{-1}X_ly)$ lies in the Lie algebra and
$d\sigma(W)(\sigma(y)u)=\sigma(y)\,d\sigma(y^{-1}Wy)u$. In particular, if $B$ is Sobolev-linear, so is
$v\mapsto B(\sigma(y)v)$.
\formallabel{GlobalTheory.T.DerivedSeriesAdjoint}
\end{lemma}

\begin{proof}
$\sigma(\exp tX)\sigma(y)u=\sigma(y)\sigma(\exp(ty^{-1}Xy))u$; differentiate at $t=0$ and induct on the length of
$W$. The second statement follows by applying the Sobolev bound of $B$ to $\sigma(y)v$, using the first statement
and $\|\sigma(y)x\|=\|x\|$.
\end{proof}

\begin{lemma}[Reproduction commutes with Sobolev-linear maps]\label{pf:sh-local:cutlim}
Let $\mu$ be a Haar measure on $\Uni(H_m)$, $w$ a smooth vector and $\kappa\in\C$ such that for every smooth $v$
the function $y\mapsto\langle v,\sigma(y)w\rangle$ is integrable and
\begin{equation}\label{eq:pf:sh-local:repro}
\int\overline{\langle v,\sigma(y)w\rangle}\,\sigma(y)w\,d\mu(y)=\kappa\,v .
\end{equation}
Let $B:\mathcal H\to V$ be Sobolev-linear with $\sup_g\|B(\sigma(g)w)\|<\infty$. Then for every smooth $v$
\[
\kappa\,B(v)=\int\overline{\langle v,\sigma(y)w\rangle}\,B(\sigma(y)w)\,d\mu(y).
\]
\formallabel{GlobalTheory.T.truncated_derivation_limit}
\end{lemma}

\begin{proof}
Write $B=\sum_i\beta_i\circ d\sigma(W_i)$ (\cref{pf:sh-local:factor}). Let $\rho$ be a smooth compactly supported
function on the matrix space with $\rho(0)=1$, and for $R>0$ and smooth $v$ put
$v_R=\int\rho(R^{-1}y)\overline{\langle v,\sigma(y)w\rangle}\sigma(y)w\,d\mu(y)$. By the truncated derivation
formula (\lref{GlobalTheory.T.TruncatedDerivation}: $v_R$ is smooth and
$d\sigma(X)v_R$ is the sum of the same integral with $\rho$ replaced by its derivative $L_X\rho$ along $X$ and the
integral with $v$ replaced by $d\sigma(X)v$) and induction, $d\sigma(W_i)v_R\to\rho(0)\kappa\,d\sigma(W_i)v$ as
$R\to\infty$, the terms containing derivatives of $\rho(R^{-1}\cdot)$ being $O(R^{-1})$; hence
$B(v_R)\to\kappa B(v)$. On the other hand $d\sigma(W_i)$ commutes with the integral defining $v_R$ (the same
formula), so $B(v_R)=\int\rho(R^{-1}y)\overline{\langle v,\sigma(y)w\rangle}B(\sigma(y)w)\,d\mu(y)$, which converges
to $\int\overline{\langle v,\sigma(y)w\rangle}B(\sigma(y)w)\,d\mu(y)$ by dominated convergence, since
$B(\sigma(y)w)$ is bounded and continuous in $y$ and $\langle v,\sigma(y)w\rangle$ is integrable. The two limits
agree.
\end{proof}

\begin{lemma}[Equivariant functions are proportional to the multiplicity-one map]\label{pf:sh-local:mfix}
Let $\mathcal D$, $X=(\sigma,w,a)$, $\mathcal V$ and $l$ be as in the multiplicity-one statement above (with
$a\ge4n^3$), and $d$ a formal degree of $\sigma$ for $\mu_G$ as in \eqref{eq:pf:sh-local:schur}. Let
$W':G\to V$ be measurable with $W'(mhx)=\shlocalpsi(m)\pi(h)W'(x)$ for $h\in U_z$, $m\in N$, $x\in G$, and such
that $y\mapsto\|x^{-1}y\|^{-a}\|W'(y)\|$ is $\mu_G$-integrable for every $x\in G$. Then there is $c\in\C$ with
\[
\int_G d\,\overline{\langle w,\sigma(y)w\rangle}\,W'(xy)\,d\mu_G(y)=c\,l(\sigma(x)w)\qquad(x\in G).
\]
\formallabel{ReciprocityTheory.ChartRestriction.DS.Mult1.stMFix}
\end{lemma}

\begin{proof}
By the folding formula (\lref{ReciprocityTheory.ChartRestriction.DS.Mult1.Fold.foldS}) there is a measurable
function $b\ge0$ on $G$ such that for every integrable $\varphi$ on $G$
\begin{equation}\label{eq:pf:sh-local:fold}
\int_G\varphi\,d\mu_G=\int_G b(y)\int_{U_z\times N}\varphi\bigl((mh)^{-1}y\bigr)\,d\mu_W(h)\,d\mu_N(m)\,d\mu_G(y).
\end{equation}
Fix $x$ and put $\varphi(y)=d\,\overline{\langle w,\sigma(x^{-1}y)w\rangle}W'(y)$; it is integrable because
$|\langle w,\sigma(x^{-1}y)w\rangle|\le C_w\|x^{-1}y\|^{-a}$
(\lref{ReciprocityTheory.ChartRestriction.DS.Mult1.MFix.phi_integrable}). By left invariance of $\mu_G$ the
left side of the claim is $\int\varphi\,d\mu_G$. In \eqref{eq:pf:sh-local:fold}, the equivariance of $W'$ gives
$W'((mh)^{-1}y)=\overline{\shlocalpsi(m)}\pi(h)^{-1}W'(y)$, and unitarity of $\sigma$ gives
$\overline{\langle w,\sigma(x^{-1}(mh)^{-1}y)w\rangle}=\langle\sigma(y)w,\sigma(mh)\sigma(x)w\rangle$; hence
\[
\varphi\bigl((mh)^{-1}y\bigr)=d\,\langle\sigma(y)w,\sigma(mh)\sigma(x)w\rangle\,\overline{\shlocalpsi(m)}\,\pi(h)^{-1}W'(y)
\]
(\lref{ReciprocityTheory.ChartRestriction.DS.Mult1.MFix.phi_Fm}), which is $d$ times the integrand of the
Bessel operator $\mathfrak B(\sigma(x)w,\sigma(y)w)$ of \eqref{eq:pf:sh-local:beop}, applied to $W'(y)$. The
vectors $\sigma(x)w$ and $\sigma(y)w$ lie in $\mathcal V$
(\lref{ReciprocityTheory.ChartRestriction.DS.Mult1.MFix.orbit_mem}), and the Bessel integral converges
absolutely for them (\lref{ReciprocityTheory.ChartRestriction.DS.Mult1.stPairs},
\lref{ReciprocityTheory.ChartRestriction.DS.Mult1.stBeConv}). So the inner integral equals
$d\,\mathfrak B(\sigma(x)w,\sigma(y)w)(W'(y))$, which by the defining property of $l$ is
$d\,\langle l(\sigma(y)w),W'(y)\rangle\,l(\sigma(x)w)$. Hence the claim holds with
$c=\int_G b(y)\,d\,\langle l(\sigma(y)w),W'(y)\rangle\,d\mu_G(y)$, which does not depend on $x$.
\end{proof}

\begin{proposition}[Orbit multiplicity]\label{pf:sh-local:orbitmult}
Let $\mathcal D$ be a diagonal real place datum for $(z,\pi)$ with $V$ finite-dimensional, $X=(\sigma,w,a)$ as
in the multiplicity-one statement with $a\ge4n^3$ and $l$ the multiplicity-one map. Assume that $\sigma$ is unitary
on $G$, $w$ is smooth, and every smooth $v$ satisfies $|\langle v,\sigma(x)w\rangle|\le C_v\|x\|^{-a}$ on $G$. Let
$B$ be Sobolev-linear with
\[
B(\sigma(hm)v)=\shlocalpsi(m)\pi(h)B(v)\qquad(h\in U_z,\ m\in N,\ v\text{ smooth}),
\]
and bounded on the orbit of $w$: $\|B(\sigma(g)w)\|\le K$ for $g\in G$. Then there is $c\in\C$ with
$B(\sigma(x)w)=c\,l(\sigma(x)w)$ for all $x\in G$.
\formallabel{CompanionTheory.Bmu.orbit_mult}
\end{proposition}

\begin{proof}
Let $d$ be the formal degree of $\sigma$ (\lref{ReciprocityTheory.ChartRestriction.DS.Mult1.stSchurG}). The
function $W'(g)=B(\sigma(g)w)$ is continuous on $G$
(\cref{pf:sh-local:factor}: it is a finite sum of bounded operators applied to $g\mapsto d\sigma(W_i)\sigma(g)w$,
which is continuous by \cref{pf:sh-local:dsad}), hence measurable. It satisfies
$W'(mhx)=\shlocalpsi(m)\pi(h)W'(x)$: indeed $mhx=h(h^{-1}mh)x$ with $h^{-1}mh\in N$ and
$\shlocalpsi(h^{-1}mh)=\shlocalpsi(m)$ for $h\in U_z$
(\lref{Formalism.QuantumDeterminants.unitary_weyl_normal_holds}), and $\sigma(x)w$ is smooth.
Moreover $y\mapsto\|x^{-1}y\|^{-a}\|W'(y)\|\le K\|x^{-1}y\|^{-a}$ is integrable, since $y\mapsto\|y\|^{-a}$ is
$\mu_G$-integrable for every $a$ at least an explicit constant $A_B(k,r)\le n^2+1$
(\lref{ReciprocityTheory.ChartRestriction.DS.Mult1.TopH.gN_integrable}; note $n^2+1\le4n^3\le a$) and
$\mu_G$ is left invariant. \Cref{pf:sh-local:mfix} gives $c$ with
\begin{equation}\label{eq:pf:sh-local:fix}
\int_G d\,\overline{\langle w,\sigma(y)w\rangle}\,B(\sigma(xy)w)\,d\mu_G(y)=c\,l(\sigma(x)w)\qquad(x\in G).
\end{equation}
On the other hand, by \eqref{eq:pf:sh-local:schur} with $\|w\|=1$, for every smooth $v$ and every $\xi$,
\[
\int\overline{\langle v,\sigma(y)w\rangle}\langle\xi,\sigma(y)w\rangle\,d\mu_G(y)=\langle\xi,v\rangle/d,
\]
the integral being absolutely convergent by the decay of $\langle v,\sigma(y)w\rangle$; hence the reproducing formula
\eqref{eq:pf:sh-local:repro} holds with $\kappa=d^{-1}$. Fix $x\in G$. The map $B_x(v)=B(\sigma(x)v)$ is
Sobolev-linear (\cref{pf:sh-local:dsad}) and bounded on the orbit of $w$, so \cref{pf:sh-local:cutlim} applied to
$B_x$ and $v=w$ gives
\[
d^{-1}B(\sigma(x)w)=\int\overline{\langle w,\sigma(y)w\rangle}\,B(\sigma(xy)w)\,d\mu_G(y).
\]
Multiplying by $d$ and comparing with \eqref{eq:pf:sh-local:fix} gives $B(\sigma(x)w)=c\,l(\sigma(x)w)$.
\end{proof}

\begin{lemma}[Absolute bound by integration along a character direction]\label{pf:sh-local:absbound}
Let $B:\mathcal H\to V$ be Sobolev-linear for a unitary representation $\sigma$ of $\Uni(H_m)$, let
$Q_1,\dots,Q_e$ be elements of the Lie algebra, and $C\ge0$, $m\in\N$ such that $\|Bv\|\le CR$ for every smooth $v$
and every $R$ with $\|d\sigma(P)v\|\le R$ for all words $P$ of length at most $m$ in the $Q_i$. Let $C_0\in\R$,
$c\in\N$ and $w$ smooth. Then there is $K$ such that $\|B(\sigma(g)w)\|\le K$ for every $g\in\Uni(H_m)$ for which
there is $M\ge1$ with $\|g^{-1}Q_ig\|\le C_0M^c$ for all $i$, and either $M=1$, or there are $Y$ in the Lie algebra and
$\lambda\in\R$ with $\|Y\|\le C_0$, $|\lambda|\ge M$ and $B(\sigma(g\exp(tY))u)=e^{i\lambda t}B(\sigma(g)u)$ for all
$t\in\R$ and smooth $u$.
\formallabel{CompanionTheory.Bmu.absBound_of}
\end{lemma}

\begin{proof}
By the word-length bound (\lref{GlobalTheory.T.word_length_bound_of_linear}) there is $C_w\ge0$ with
$\|d\sigma(W)w\|\le C_w\prod_j\|X_j\|$ for all words $W=(X_1,\dots,X_l)$ in the Lie algebra of length
$l\le m+cm$. Put $A=\max(1,C_0)$ and $K=C\,C_wA^{m+cm}$.

\emph{Claim}: for every word $R$ of length at most $cm$ whose letters have norm at most $A$,
$\|B(\sigma(g)d\sigma(R)w)\|\le K\,M^{cm}$. Indeed, for a word $P$ of length at most $m$ in the $Q_i$,
\cref{pf:sh-local:dsad} gives $\|d\sigma(P)\sigma(g)d\sigma(R)w\|=\|d\sigma(g^{-1}Pg)d\sigma(R)w\|$, and
$d\sigma(g^{-1}Pg)d\sigma(R)=d\sigma(g^{-1}Pg\,R)$ is the derivative along the concatenated word, of length at most
$m+cm$; its letters $g^{-1}Q_ig$ have norm at most $C_0M^c\le AM^c$. Hence
$\|d\sigma(P)\sigma(g)d\sigma(R)w\|\le C_w(AM^c)^mA^{cm}=C_wA^{m+cm}M^{cm}$, and the hypothesis on $B$ gives the claim.

If $M=1$, the claim with the empty word gives $\|B(\sigma(g)w)\|\le K$. Otherwise let $Y,\lambda$ be as in the
hypothesis. For smooth $u$, differentiate $t\mapsto B(\sigma(g\exp(tY))u)=e^{i\lambda t}B(\sigma(g)u)$ at $t=0$:
writing $B=\sum_i\beta_i\circ d\sigma(W_i)$ (\cref{pf:sh-local:factor}) and using the derivative formula
$\frac d{dt}d\sigma(W_i)\sigma(g\exp(tY))u|_{t=0}=d\sigma(W_i)\sigma(g)d\sigma(Y)u$
(\lref{GlobalTheory.T.DerivedSeriesDerivation}), we get $B(\sigma(g)d\sigma(Y)u)=i\lambda\,B(\sigma(g)u)$. By
induction, $B(\sigma(g)d\sigma(Y^{j})w)=(i\lambda)^jB(\sigma(g)w)$ for the word $Y^j=(Y,\dots,Y)$. The claim with
$R=Y^{cm}$ (note $\|Y\|\le C_0\le A$) gives $|\lambda|^{cm}\|B(\sigma(g)w)\|\le KM^{cm}$, and $|\lambda|\ge M>0$ gives
$\|B(\sigma(g)w)\|\le K$.
\end{proof}

\begin{theorem}[Bessel maps are bounded on orbits]\label{pf:sh-local:bounded}
Let $\sigma$ be a unitary representation of $G=\Uni(H_{k,r})$, $k\ge1$, $z\in W_0$ with $T(z)>0$,
$c\in(\C^\times)^r$, $V$ a finite-dimensional Hilbert space and $\pi$ a map from $U_z$ to the operators on $V$ with
$\pi(1)=1$, all $\pi(h)$ isometric and $\pi$ continuous on $U_z$. Let $B$ be Sobolev-linear with
$B(\sigma(hm)v)=\psi_c(m)\pi(h)B(v)$ for $h\in U_z$, $m\in N$ and smooth $v$. Then for every smooth $w$ there is
$K$ with $\|B(\sigma(g)w)\|\le K$ for all $g\in G$.
\formallabel{CompanionTheory.Bmu.bounded}
\formallabel{CompanionTheory.Bmu.bounded_of}
\end{theorem}

\begin{proof}
The proof combines eight structural facts about $G$, whose formal proofs we do not reproduce
(\lref{GlobalTheory.T.AdaptedBasisExistsAlt}, \lref{GlobalTheory.T.IwasawaDecomposition}, \lref{GlobalTheory.T.LeviDecomposition}, \lref{GlobalTheory.T.CartanDecomposition}, \lref{GlobalTheory.T.SpanGenerators}, \lref{GlobalTheory.T.ExponentialMembership}, \lref{GlobalTheory.T.RootVectors}, and \lref{GlobalTheory.T.torBound'})
with the quotient bound (\lref{GlobalTheory.T.quotient_bound_of_factor_reduce_chart}) and
\cref{pf:sh-local:absbound}.
\begin{enumerate}[label=(\alph*)]
\item \emph{Adapted basis.} There is $\Gamma\in\GL_n(\C)$ fixing all $e_a,e_a^*$ with $\Gamma^*H\Gamma=H_q$, where
$H_q$ is the antidiagonal form on the outer coordinates $i\le r$ or $i\ge r+k$ and the identity on the inner ones,
and $z=\rho\,\Gamma(e_r-e_{r+k})$ for some $\rho>0$.
\item \emph{Decomposition} (from the Iwasawa, Levi and Cartan decompositions of $\Uni(H_q)$,
\lref{GlobalTheory.T.decomposition_of_iwasawa_levi_cartan}). Every $g\in G$ can be written
$g=h\,m\,\Gamma\,\delta\,R$ with $h\in U_z$, $m\in N$, $R$ unitary, $\Gamma\delta R\in G$, and $\delta$ a positive
diagonal matrix of \emph{torus type}: $\delta_i\delta_j=1$ for outer $i,j$ with $i+j=n-1$, $\delta_i=1$ for inner
$i$, and $\delta_r\le1$.
\item \emph{Spanning.} $\mathfrak g=\mathfrak u_z+\mathfrak n+\mathfrak q_\Gamma$, where $\mathfrak u_z$ and $\mathfrak n$
are the Lie algebras of $U_z$ and $N$ and $\mathfrak q_\Gamma$ consists of the $Y\in\mathfrak g$ with
$(\Gamma^{-1}Y\Gamma)_{ij}=0$ for $i<j$ with $i$ or $j$ outer. Exponentials of elements of $\mathfrak u_z$, resp.
$\mathfrak n$, lie in $U_z$, resp. $N$.
\item \emph{Quotient bound.} Choose bases $X_j$ of $\mathfrak u_z$ and $\mathfrak n$ and $Q_i$ of $\mathfrak q_\Gamma$.
Along $X_j$ the map $B$ transforms by $\pi$, resp. by $\psi_c$, i.e. $B(\sigma(\exp tX_j)v)=\chi_j(t)B(v)$ with
continuous $\chi_j$, $\chi_j(0)=1$. Using a local chart of $G$ adapted to the decomposition (c)
(\lref{GlobalTheory.T.ChartExistence}) and the reduction of derivatives along words
(\lref{GlobalTheory.T.reduction_of_derived_series}), one obtains $C\ge0$ and $m\in\N$ with
$\|Bv\|\le C\max_{|P|\le m}\|d\sigma(Q_P)v\|$ for smooth $v$, $Q_P$ the words in the $Q_i$.
\item \emph{Root vectors.} There is $C_1$ such that for every $a<r$ and every $\delta$ of torus type with
$g'=\Gamma\delta R\in G$ there is $Y\in\mathfrak g$ with $\|Y\|\le C_1$, $g'\exp(tY)g'^{-1}\in N$ and
$\psi_c(g'\exp(tY)g'^{-1})=\exp(i\,t\,\delta_a/\delta_{a+1})$ for all $t$.
\item \emph{Torus bound.} If $\delta$ is of torus type, $R$ unitary, $M\ge1$ and $\delta_a\le M\delta_{a+1}$ for
all $a<r$, then $\|(\delta R)^{-1}A(\delta R)\|\le n^2\|A\|M^n$ for every $A$ with $A_{ij}=0$ for $i<j$, $i$ or $j$
outer.
\item \emph{Absolute bound}: \cref{pf:sh-local:absbound}, applied with the quotient bound (d).
\end{enumerate}
Now let $g\in G$ and write $g=hm g'$, $g'=\Gamma\delta R$, by (b). Since $\sigma(g')w$ is smooth,
$\|B(\sigma(g)w)\|=|\psi_c(m)|\,\|\pi(h)B(\sigma(g')w)\|=\|B(\sigma(g')w)\|$. Let
$M=\max(1,\max_a\delta_a/\delta_{a+1})$. By (f), $\|g'^{-1}Q_ig'\|\le C_2M^n$ with
$C_2=n^2\sum_i\|\Gamma^{-1}Q_i\Gamma\|$. If $M=1$, (g) applies with $C_0=\max(C_1,C_2)$. Otherwise let $a_0$
maximise $\delta_a/\delta_{a+1}$ and take $Y$ from (e) for $a_0$; then
$g'\exp(tY)=(g'\exp(tY)g'^{-1})g'$ with $g'\exp(tY)g'^{-1}\in N$, so by the equivariance of $B$,
$B(\sigma(g'\exp(tY))u)=e^{i\lambda t}B(\sigma(g')u)$ with $\lambda=\delta_{a_0}/\delta_{a_0+1}=M$, and (g) applies
again. In both cases $\|B(\sigma(g)w)\|\le K$ with $K$ independent of $g$.
\end{proof}

\subsubsection{Proof of \texorpdfstring{\cref{pf:sh-local:bmu}}{the uniqueness theorem}}

\begin{proof}[Proof of \cref{pf:sh-local:bmu}]
The proof combines \cref{pf:sh-local:bounded}, \cref{pf:sh-local:cutlim}, \cref{pf:sh-local:torconj},
\cref{pf:sh-local:dsad} and \cref{pf:sh-local:factor} (through \cref{pf:sh-local:orbitmult}).
Let the data be as in the theorem; $\mathcal H$ is complete and $V$ is finite-dimensional, $\pi$ is a continuous
unitary irreducible representation of $U_z$, and there is a diagonal real place datum $\mathcal D$ for $(z,\pi)$
(see the inputs above). By \lref{ReciprocityTheory.ChartRestriction.DS.Mult1.TopTA.sqInt_of_DSHpk} $\sigma$ is
unitary on $G$ with square integrable coefficients.

\emph{Decay.} Since $\theta^*(\eps)\le40n^3$ and $h\ge120n^3$, we have $\theta^*(\eps)+4n^3<h$; so the decay
statement with $m=4n^3$ gives, for all smooth $v,w'$, a bound
$|\langle v,\sigma(x)w'\rangle|\le C\|x\|^{-4n^3}$ on $G$. Let $w_0$ be a $K$-finite unit vector of the lowest
$K$-type; it is smooth. With $w=w_0$, $a=4n^3$ and $C_w$ the constant of the decay bound for $v=w'=w_0$,
$X=(\sigma,w_0,a)$ satisfies the hypotheses of the multiplicity-one statement; let $l$ be the multiplicity-one map
and $\tau$ the element of \cref{pf:sh-local:torconj}.

\emph{Each Bessel map is a multiple of $l$ on the orbit of $w_0$.} Let $B$ be a smooth Bessel map for
$(z,\psi_c,\sigma,\pi)$. It is Sobolev-linear, and by \cref{pf:sh-local:bounded} there is $K$ with
$\|B(\sigma(g)w_0)\|\le K$ for $g\in G$. Put $B'(v)=B(\sigma(\tau)v)$. It is Sobolev-linear
(\cref{pf:sh-local:dsad}), bounded by $K$ on the orbit of $w_0$, and for $h\in U_z$, $m\in N$ and smooth $v$,
writing $\tau hm=h(\tau m\tau^{-1})\tau$,
\[
B'(\sigma(hm)v)=B(\sigma(h\,\tau m\tau^{-1})\sigma(\tau)v)=\psi_c(\tau m\tau^{-1})\pi(h)B(\sigma(\tau)v)=\shlocalpsi(m)\pi(h)B'(v).
\]
By \cref{pf:sh-local:orbitmult} there is $c_B$ with $B(\sigma(\tau)\sigma(x)w_0)=c_B\,l(\sigma(x)w_0)$ for all
$x\in G$.

\emph{Conclusion.} Let $c_1,c_2$ be these constants for $B_1,B_2$. If $c_1=0$ take $(a,b)=(1,0)$, otherwise
$(a,b)=(c_2,-c_1)$; in both cases $(a,b)\ne0$ and $ac_1+bc_2=0$. The map $B''=aB_1+bB_2$ is Sobolev-linear and
$B''(\sigma(\tau)\sigma(x)w_0)=(ac_1+bc_2)\,l(\sigma(x)w_0)=0$ for all $x$. Let $v$ be smooth and
$v'=\sigma(\tau^{-1})v$, which is smooth. The map $u\mapsto B''(\sigma(\tau)u)$ is Sobolev-linear
(\cref{pf:sh-local:dsad}) and vanishes on all $\sigma(x)w_0$, hence on all finite linear combinations
$\sum_ic_i\sigma(g_i)w_0$. By cyclicity (\lref{CompanionTheory.D.Upq.Sm.cyc_of_smooth}) these combinations
approximate $v'$ in every finite family of Sobolev seminorms, so the Sobolev bound gives $B''(\sigma(\tau)v')=0$,
i.e. $aB_1(v)+bB_2(v)=0$.
\end{proof}

% ---------------------------------------------------------------------------
\subsection{Real objects at infinity}\label{sec:pf:sh-local:real}

The energy bound for the real objects of a Hida setting is an instance of an abstract inequality, which we state
and prove first.

\subsubsection{An abstract energy bound}
Let $G$ be a second countable, locally compact Hausdorff group with a left and right invariant Haar measure $\mu$.

\begin{definition}[Haar subgroups, admissible data, disintegrations]\label{pf:sh-local:def:wabs}
For a subset $S\subseteq G$ and a measure $\mu_S$ on $G$ we say that $(S,\mu_S)$ is a \emph{Haar subgroup} if $S$ is a
closed subgroup and $\mu_S$ is $\sigma$-finite, non-zero, concentrated on $S$, finite on compact sets, invariant
under left and right translations by elements of $S$ and under inversion; it is \emph{F\o lner} if for every compact
$K\subseteq S$ and $\eta>0$ there is a measurable $\Omega\subseteq S$ with $0<\mu_S(\Omega)<\infty$ and
$\mu_S(\Omega)\le(1+\eta)\,\mu_S\{t\in\Omega:Kt\subseteq\Omega\}$. A measure $\nu$ is a \emph{quotient measure} for
$(\mu,\mu_S)$ if $\int f\,d\mu=\int\!\int f(sx)\,d\mu_S(s)\,d\nu(x)$ for every measurable $f\ge0$.

Let $\sigma$ be a unitary representation of $G$ on a Hilbert space $\mathcal H$ (with $\sigma(1)=1$,
$\sigma(xy)=\sigma(x)\sigma(y)$ and continuous orbit maps), $d>0$ and $w\in\mathcal H$. The pair $(\sigma,d)$
satisfies \emph{Schur orthogonality} if all functions $x\mapsto\langle b,\sigma(x)a\rangle\overline{\langle
b',\sigma(x)a'\rangle}$ are $\mu$-integrable with integral $\langle a',a\rangle\langle b,b'\rangle/d$. Given a
Haar subgroup $(S,\mu_S)$, a quotient measure $\nu$, a finite-dimensional Hilbert space $V$, a map $\tau$ from $S$
to the unitary operators of $V$ with $\tau(1)=1$ and $\tau(st)=\tau(s)\tau(t)$, and a continuous $W:G\to V$, the
datum $(W,\tau,\nu)$ is \emph{admissible} (for $\sigma,d,w,S$) if $W(sx)=\tau(s)W(x)$ for $s\in S$, the
\emph{reproducing identity}
\begin{equation}\label{eq:pf:sh-local:wrepro}
\int_G d\,\overline{\langle w,\sigma(y)w\rangle}\,W(xy)\,d\mu(y)=W(x)\qquad(x\in G)
\end{equation}
holds with an absolutely convergent integral, $\int\|W\|\,d\nu<\infty$ and $\int\|W\|^2\,d\nu=1$. Finally, a measure
space $(C,m)$ with a map $h:C\to G$ is a \emph{disintegration} of $\mu$ along $(\mu_S,\mu_{S'})$ if
\begin{equation}\label{eq:pf:sh-local:disint}
\int_G\varphi\,d\mu=\int_C\int\!\!\int\varphi\bigl(s\,h(c)\,s'^{-1}\bigr)\,d\mu_{S'}(s')\,d\mu_S(s)\,dm(c)
\end{equation}
for every measurable $\varphi\ge0$.
\formal{Formalism.L2.SG.Pf.EnergyAbs.IsWAbs}\formal{Formalism.L2.SG.Pf.EnergyAbs.IsDisint}
\end{definition}

\begin{theorem}[Abstract energy bound]\label{pf:sh-local:energyabs}
Let $\sigma$ be a unitary representation of $G$ satisfying Schur orthogonality with $d>0$, and $w$ a unit vector.
Let $(S,\mu_S)$ and $(S',\mu_{S'})$ be F\o lner Haar subgroups with $s$-finite quotient measures $\nu,\nu'$, let
$(W,\tau,\nu)$ and $(W',\tau',\nu')$ be admissible, with values in $V$ and $V'$, and let $(C,m,h)$ be a
disintegration of $\mu$ along $(\mu_S,\mu_{S'})$ with $h$ measurable. Assume:
\begin{enumerate}[label=(\roman*)]
\item for every $c\in C$ the function $x\mapsto W(x)\otimes W'(h(c)^{-1}x)^*$ is $\mu$-integrable;
\item for every $c\in C$,
\[
\int\!\!\int\!\!\int\!\!\int|\langle\sigma(z)w,\sigma(s\,h(c)\,s'^{-1})\sigma(z')w\rangle|\,\|W(z)\|\,\|W'(z')\|
\,d\mu_{S'}(s')\,d\mu_S(s)\,d\nu'(z')\,d\nu(z)<\infty;
\]
\item the integral $\int\!\int\!\int|\langle\sigma(z)w,\sigma(s)\sigma(z')w\rangle|\,\|W(z)\|\,\|W(z')\|\,d\mu_S(s)\,
d\nu(z')\,d\nu(z)$ is finite, and so is the analogous integral for $(S',W',\nu')$.
\end{enumerate}
Then
\[
\int_C\shlocalHS{\int_G W(x)\otimes W'(h(c)^{-1}x)^*\,d\mu(x)}^2\,dm(c)\ \le\ \frac1d .
\]
\formal{Formalism.L2.SG.Pf.EnergyAbs.EnergyAbs}
\formallabel{Formalism.L2.S.Pv.EnergyAbs.energyAbs_holds}
\end{theorem}

For the proof, fix measurable sets $\Omega\subseteq S$, $\Omega'\subseteq S'$ of finite positive measure. For $e\in V$ let
\[
y_\Omega(e)=\int\!\int_{\Omega}\langle W(s^{-1}z),e\rangle\,\sigma(s^{-1}z)w\,d\mu_S(s)\,d\nu(z)\in\mathcal H,
\]
let $Y_\Omega:V\to\mathcal H$ be the linear map with $Y_\Omega e_i=y_\Omega(e_i)$ for an orthonormal basis $(e_i)$ of
$V$, define $Y'_{\Omega'}:V'\to\mathcal H$ in the same way from $(W',\nu',\mu_{S'})$, and put
$X_{\Omega,\Omega'}(g)=d\,Y_\Omega^*\,\sigma(g)\,Y'_{\Omega'}:V'\to V$ for $g\in G$.

\begin{lemma}\label{pf:sh-local:energyX}
If $(\sigma,d)$ satisfies Schur orthogonality, then $g\mapsto\shlocalHS{X_{\Omega,\Omega'}(g)}^2$ is continuous and
\[
\int_G\shlocalHS{X_{\Omega,\Omega'}(g)}^2\,d\mu(g)=d\,\shlocalHS{Y_\Omega}^2\,\shlocalHS{Y'_{\Omega'}}^2 .
\]
\formallabel{Formalism.L2.S.Pv.EnergyAbs.lintegral_hsA_Xop}
\end{lemma}

\begin{proof}
Continuity follows from the continuity of the orbit maps. With $y_i=Y_\Omega e_i$ and $y'_j=Y'_{\Omega'}e'_j$,
$\shlocalHS{X_{\Omega,\Omega'}(g)}^2=d^2\sum_{i,j}|\langle\sigma(g)y'_j,y_i\rangle|^2$, and by Schur orthogonality
$\int|\langle\sigma(g)y'_j,y_i\rangle|^2\,d\mu=\|y'_j\|^2\|y_i\|^2/d$. Summing gives
$d\sum_i\|y_i\|^2\sum_j\|y'_j\|^2$.
\end{proof}

\begin{lemma}[Normalisation along F\o lner sets]\label{pf:sh-local:folnerlimit}
Let $(\sigma,d)$ satisfy Schur orthogonality, $\|w\|=1$, $(S,\mu_S)$ a Haar subgroup with an $s$-finite quotient
measure $\nu$, $(W,\tau,\nu)$ admissible, and assume condition (iii) of \cref{pf:sh-local:energyabs} for
$(S,W,\nu)$. Let $\Omega_n\subseteq S$ be measurable with $0<\mu_S(\Omega_n)<\infty$ and
$r_n(q):=\mu_S\{t\in\Omega_n:qt\in\Omega_n\}/\mu_S(\Omega_n)\to1$ for every $q\in S$. Then
$\shlocalHS{Y_{\Omega_n}}^2/\mu_S(\Omega_n)\to1/d$.
\formallabel{Formalism.L2.S.Pv.EnergyAbs.folner_limit}
\end{lemma}

\begin{proof}
Let $F(z,z',q)=\langle W(qz'),W(z)\rangle\,\langle\sigma(z)w,\sigma(qz')w\rangle$ on $G\times G\times G$.
(a) $F$ is integrable for $\nu\times\nu\times\mu_S$ (\lref{Formalism.L2.S.Pv.EnergyAbs.Fk_integrable}): for
$q\in S$, $\|W(qz')\|=\|\tau(q)W(z')\|=\|W(z')\|$, so $|F|$ is bounded by the integrand of (iii), and $\mu_S$ is
concentrated on $S$. (b) For every measurable $\Omega\subseteq S$ of finite positive measure,
$\shlocalHS{Y_\Omega}^2/\mu_S(\Omega)=\int r_\Omega(q)F(z,z',q)\,d(\nu\times\nu\times\mu_S)$ with
$r_\Omega(q)=\mu_S\{t\in\Omega:qt\in\Omega\}/\mu_S(\Omega)$ (\lref{Formalism.L2.S.Pv.EnergyAbs.sumsq_div_eq}, a
direct computation from the definition of $y_\Omega$, using the equivariance of $W$ under $S$ and the Haar subgroup
properties of $\mu_S$). (c) $\int F\,d(\nu\times\nu\times\mu_S)=1/d$ (\lref{Formalism.L2.S.Pv.EnergyAbs.Fk_value}): since
$\mu$ is the image of $\nu\times\mu_S$ under $(z',q)\mapsto qz'$, the integral equals
\begin{align*}
\int\!\!\int\langle W(x),W(z)\rangle\langle\sigma(z)w,\sigma(x)w\rangle\,d\mu(x)\,d\nu(z)
&=\int\Bigl\langle\int\overline{\langle w,\sigma(y)w\rangle}\,W(zy)\,d\mu(y),\,W(z)\Bigr\rangle d\nu(z)\\
&=\int\frac{\|W(z)\|^2}d\,d\nu(z)=\frac1d,
\end{align*}
by the substitution $x=zy$, unitarity, the reproducing identity \eqref{eq:pf:sh-local:wrepro} and
$\int\|W\|^2\,d\nu=1$. Since $0\le r_n\le1$ and $r_n(q)\to1$ for $q\in S$, dominated convergence and (a)--(c) give the
claim.
\end{proof}

\begin{proof}[Proof of \cref{pf:sh-local:energyabs}]
(1) \emph{F\o lner sequences.} Since $G$ is second countable and $(S,\mu_S)$ is a F\o lner Haar subgroup, there are
measurable $\Omega_n\subseteq S$ with $0<\mu_S(\Omega_n)<\infty$ such that
$\mu_S\{t\in\Omega_n:Kt\subseteq\Omega_n\}/\mu_S(\Omega_n)\to1$ for every compact $K\subseteq S$
(\lref{Formalism.L2.S.Pv.EnergyAbs.folner_seq}); similarly $\Omega'_n\subseteq S'$. Write $Y_n,Y'_n,X_n$ for
$Y_{\Omega_n},Y'_{\Omega'_n},X_{\Omega_n,\Omega'_n}$.

(2) \emph{Normalisation.} By \cref{pf:sh-local:folnerlimit} (applied with $K=\{q\}$ in the F\o lner property),
\[
\shlocalHS{Y_n}^2/\mu_S(\Omega_n)\to1/d\quad\text{and}\quad\shlocalHS{Y'_n}^2/\mu_{S'}(\Omega'_n)\to1/d .
\]

(3) \emph{Pointwise lower bound.} For every $\gamma\in G$ satisfying the convergence condition (ii) with $h(c)$
replaced by $\gamma$,
\begin{align*}
&\shlocalHS{\int_GW(x)\otimes W'(\gamma^{-1}x)^*\,d\mu(x)}^2\\
&\qquad\le\liminf_{n\to\infty}\frac1{\mu_S(\Omega_n)\mu_{S'}(\Omega'_n)}
\int\!\!\int\shlocalHS{X_n(s\gamma s'^{-1})}^2\,d\mu_{S'}(s')\,d\mu_S(s)
\end{align*}
(\lref{Formalism.L2.S.Pv.EnergyAbs.point_lower}).

(4) \emph{Integration.} Let $F_n(c)$ be the right side of (3) without the $\liminf$, for $\gamma=h(c)$; it is
measurable in $c$ because $h$ is measurable and $g\mapsto\shlocalHS{X_n(g)}^2$ is continuous. By (3) and Fatou's
lemma,
\[
\int_C\shlocalHS{\int_G W(x)\otimes W'(h(c)^{-1}x)^*\,d\mu}^2\,dm(c)\le\int_C\liminf_nF_n\,dm\le\liminf_n\int_CF_n\,dm .
\]
By the disintegration \eqref{eq:pf:sh-local:disint} applied to $\varphi=\shlocalHS{X_n(\cdot)}^2$ and
\cref{pf:sh-local:energyX},
\[
\int_CF_n\,dm=\frac{d\,\shlocalHS{Y_n}^2\,\shlocalHS{Y'_n}^2}{\mu_S(\Omega_n)\mu_{S'}(\Omega'_n)}
=d\cdot\frac{\shlocalHS{Y_n}^2}{\mu_S(\Omega_n)}\cdot\frac{\shlocalHS{Y'_n}^2}{\mu_{S'}(\Omega'_n)}\longrightarrow
d\cdot\frac1d\cdot\frac1d=\frac1d
\]
by (2) (\lref{Formalism.L2.S.Pv.EnergyAbs.tendsto_energy}). Hence the $\liminf$ is $1/d$.
\end{proof}

\subsubsection{The real objects of a Hida setting}
Let $D$ be a datum and $X$ a Hida setting of $D$ (\cref{thm:theorem-S}). We use the following parts of $X$. It contains an automorphic
instance: a totally real field $F$ with set $\Sigma_\infty$ of infinite places, a CM extension $E/F$, a hermitian
space of dimension $n=r$ (the rank of $D$) with signatures $\eps_v$ ($v\in\Sigma_\infty$), and the real group
$G_\infty=\prod_v\Uni(J_{\eps_v})$ with the product $\mu_\infty$ of Haar measures. At every $v$ it contains a discrete
series $\sigma_v$ of $\Uni(J_{\eps_v})$ with a unit vector $w_v$ of the lowest $K$-type and formal degree $d_v$; let
$\sigma_\otimes$ be the tensor product representation of $G_\infty$, $w_\otimes=\bigotimes_vw_v$ and
$\deg X=\prod_vd_v$. For every word $u$ it contains: a frame $(e_a)_{a<r_S}$ (independent of $u$) and a vector
$z_u$; closed subgroups $N_\infty(u)\subseteq S_\infty(u)$ and $H_\infty(u)$ of $G_\infty$ (unipotent on the flag of
the $e_a$ and $z_u$, resp.\ fixing the frame, its dual frame and $z_u$) with Haar measures $\mu_N(u)$, $\mu_H(u)$;
the image $\mu_S(u)$ of $\mu_N(u)\times\mu_H(u)$ under multiplication and a quotient measure $\nu_u$ for
$(\mu_\infty,\mu_S(u))$; a finite-dimensional Hilbert space $V_u$ with a unitary representation $\pi_u$ of
$S_\infty(u)$, trivial on $N_\infty(u)$, and a character $\psi_{\infty,u}$; and a continuous function
$W_u:G_\infty\to V_u$ with $W_u(sx)=\psi_{\infty,u}(s)\pi_u(s)W_u(x)$ for $s\in S_\infty(u)$, reproduced by
$\varphi_\infty=\deg X\cdot\overline{\langle w_\otimes,\sigma_\otimes(\cdot)w_\otimes\rangle}$ as in
\eqref{eq:pf:sh-local:wrepro}, with $\int\|W_u\|\,d\nu_u<\infty$ and $\int\|W_u\|^2\,d\nu_u=1$.

Fix words $u,u'$ and a \emph{scale} $s_c>0$. For $v\in\Sigma_\infty$ let $f_{u,v}$ be the family of $r_S+1$
vectors of $\C^n$ formed by $s_c^{-(r_S-a)}e_{a,v}$ ($a<r_S$) followed by $z_{u,v}$, the index $v$ denoting the
coordinates at $v$. For $\gamma\in G_\infty$ let $Z_v(\gamma)$ be the $(r_S+1)\times(r_S+1)$ matrix with entries
$f_{u,v,a}^*J_{\eps_v}\gamma_vf_{u',v,b}$. For a square matrix $Z$ of size $r_S+1$ let $D_l(Z)$ be its leading
principal minor of size $l$ and
\[
\mathcal G(Z)=|D_{r_S+1}(Z)|^2-T_uT_{u'}|D_{r_S}(Z)|^2,\qquad \mathcal J(Z)=\prod_{l=1}^{r_S}|D_l(Z)|^4,
\]
where $T_u$ is the integer $T$ attached to the word $u$ by $D$. Let $C_\infty=\prod_{v}\C^{r_S+1}$ with the product
$dc$ of the measures $2\,d\mathrm{Leb}$ on $\C$, and
\[
\mathrm{Good}(u,u')=\{c\in C_\infty:\ c_{v,l}\ne0\text{ and }\mathcal G(\mathrm{diag}(c_v))>0\text{ for all }v,l\},
\]
\[
\mathrm{wt}(c)=\prod_v\mathcal G(\mathrm{diag}(c_v))^{k-1}\mathcal J(\mathrm{diag}(c_v)),
\]
$k$ the rank parameter of $D$. Say that $\gamma$ \emph{diagonalises} $c$ if $Z_v(\gamma)=\mathrm{diag}(c_v)$ for all
$v$, and let $h(c)$ be a chosen diagonalising element if one exists, and $1$ otherwise. For $C>0$ the
\emph{disintegration identity} $\mathrm{Dis}(C)$ is \eqref{eq:pf:sh-local:disint} with $\mu=\mu_\infty$,
$\mu_S=\mu_S(u)$, $\mu_{S'}=\mu_S(u')$ and $dm=C\,\mathrm{wt}(c)\,dc$ restricted to $\mathrm{Good}(u,u')$; let
$C(u,u',s_c)$ be a chosen $C>0$ with $\mathrm{Dis}(C)$ if one exists, and $0$ otherwise, and
$\nu_M=C(u,u',s_c)\,\mathrm{wt}(c)\,dc$ on $\mathrm{Good}(u,u')$. The \emph{orbital operator at infinity} is
\[
O_\infty(\gamma)=\int_{G_\infty}W_u(x)\otimes W_{u'}(\gamma^{-1}x)^*\,d\mu_\infty(x):\ V_{u'}\to V_u .
\]

\begin{definition}[Real objects at infinity]\label{pf:sh-local:def:ininf}
$\mathrm{InInf}(X,u,u',s_c)$ is the conjunction of: (i) $\mathrm{Dis}(C)$ holds for some $C>0$; (ii) every
$c\in\mathrm{Good}(u,u')$ has a diagonalising element, and $x\mapsto W_u(x)\otimes W_{u'}(h(c)^{-1}x)^*$ is
$\mu_\infty$-integrable; (iii) the \emph{energy} satisfies
\[
\shlocalE(u,u',s_c):=\int_{\mathrm{Good}(u,u')}\shlocalHS{O_\infty(h(c))}^2\,d\nu_M(c)\ \le\ \frac1{\deg X}.
\]
\formal{Formalism.L2.SG.RealObj.InInfinityX}
\end{definition}

\begin{definition}[Gap hypothesis]\label{pf:sh-local:def:gap}
The \emph{gap hypothesis} for $D$ is: $r\ge1801$ and, for every type $\tau$ and every $j$, the Harish-Chandra
parameter $b$ of $D$ at $\tau$ satisfies $b_{j+1}-b_j\ge16r^{19}$.
\formal{Formalism.L2.SG.ConcreteSettingExists.GapHyp}
\end{definition}

\begin{theorem}[Real objects at infinity]\label{pf:sh-local:realobj}
Assume $\sfname{UniqGap}$ (\cref{pf:sh-local:uniqgap}). For every datum $D$ satisfying the gap hypothesis, every Hida
setting $X$ of $D$, all words $u,u'$ and every $s_c>0$, $\mathrm{InInf}(X,u,u',s_c)$ holds.
\formallabel{Formalism.L2.SG.RealObj.EnergyConvJ.BM.realObj_inInfinityX_jsz}
\formallabel{Formalism.L2.SG.Pf.realObj_inInfinityX_ofG}
\formallabel{Formalism.L2.SG.RealObj.EnergyInst.inInfinityXAt}
\end{theorem}

The hypothesis $\sfname{UniqGap}$ is a theorem (\cref{pf:sh-local:uniqgap}); the citing sections combine the two.
The proof uses two convergence statements deduced from $\sfname{UniqGap}$, together with the decay of matrix
coefficients at the lowest $K$-type and at infinity and the smoothness of $K$-finite vectors
(\lref{CompanionTheory.Lime.BM.inputs2_of}, which packages these four inputs):
\begin{itemize}
\item $(\mathrm{Int})$ (\lref{Formalism.L2.SG.RealObj.EnergyConvJ.BM.stepIntG_jsz}): under the gap
hypothesis, for every $c\in\mathrm{Good}(u,u')$ and every $\gamma$ diagonalising $c$, the function
$x\mapsto W_u(x)\otimes W_{u'}(\gamma^{-1}x)^*$ is $\mu_\infty$-integrable;
\item $(\mathrm{Conv})$ (\lref{Formalism.L2.SG.RealObj.EnergyConvJ.BM.stepConvG_jsz}): under the gap hypothesis,
condition (ii) of \cref{pf:sh-local:energyabs} holds for the data $\sigma_\otimes$, $w_\otimes$, $W_u$, $W_{u'}$,
$\nu_u$, $\nu_{u'}$, $\mu_S(u)$, $\mu_S(u')$ at every $\gamma$ diagonalising some $c\in\mathrm{Good}(u,u')$, and
condition (iii) holds for every word.
\end{itemize}

\begin{proposition}[The energy step]\label{pf:sh-local:energystep}
Let $D$ be a datum such that $(\mathrm{Int})$ and $(\mathrm{Conv})$ hold for all Hida settings of $D$. Then for every
Hida setting $X$, words $u,u'$ and $s_c>0$: if $\mathrm{Dis}(C)$ holds for some $C>0$ and condition (ii) of
\cref{pf:sh-local:def:ininf} holds, then $\shlocalE(u,u',s_c)\le1/\deg X$.
\formallabel{Formalism.L2.SG.RealObj.EnergyInst.stepEnergyAt_of_inst}
\end{proposition}

\begin{proof}
We apply \cref{pf:sh-local:energyabs} with $G=G_\infty$, $\mu=\mu_\infty$, $\sigma=\sigma_\otimes$, $w=w_\otimes$,
$d=\deg X$, $(S,\mu_S,\nu,W)=(S_\infty(u),\mu_S(u),\nu_u,W_u)$, the same for $u'$,
$\tau(s)=\psi_{\infty,u}(s)\pi_u(s)$, $(C,m)=(C_\infty,\nu_M)$, and a measurable modification $g$ of $h$. Its
hypotheses hold:
\begin{itemize}
\item $\sigma_\otimes$ is unitary and satisfies Schur orthogonality with $d=\deg X$, $\|w_\otimes\|=1$ and
$\varphi_\infty=\deg X\,\overline{\langle w_\otimes,\sigma_\otimes(\cdot)w_\otimes\rangle}$
(\lref{Formalism.L2.S.Pv.SchurTens.stepSchurG_holds}, proved from Schur orthogonality of the $\sigma_v$ by a
density argument for tensor products);
\item $(S_\infty(u),\mu_S(u))$ is a F\o lner Haar subgroup, $\tau$ is multiplicative on $S_\infty(u)$ and $\nu_u$
is $s$-finite (\lref{Formalism.L2.S.Pv.HaarSubS.stepSubG_holds}); $\nu_u$ is a quotient measure by the
definition of a Hida setting; and $(W_u,\tau,\nu_u)$ is admissible
(\lref{Formalism.L2.SG.RealObj.EnergyInst.isWAbs_of}: the conditions are those of the Hida setting, rewritten
with the two previous items);
\item \emph{a measurable section}: there is a measurable $h_0:C_\infty\to G_\infty$ such that $h_0(c)$ diagonalises
$c$ for every $c\in\mathrm{Good}(u,u')$ (\lref{Formalism.L2.S.Pv.SecW.stepSection_holds}). If $\mathrm{Good}$
is empty, $\nu_M=0$ and there is nothing to prove; otherwise fix $c_0\in\mathrm{Good}$ and let $g=h_0$ on
$\mathrm{Good}$ and $g=h_0(c_0)$ outside; $g$ is measurable and $g(c)$ diagonalises some element of $\mathrm{Good}$
for every $c$;
\item \emph{disintegration along $g$}: two elements diagonalising the same $c$ differ by left and right
multiplication by elements of $S_\infty(u)$ and $S_\infty(u')$
(\lref{Formalism.L2.SG.RealObj.InfCut.UnitaryIsometries.StepDiagDC_holds}); by the invariance of
$\mu_S(u),\mu_S(u')$ (\lref{Formalism.L2.SG.RealObj.EnergyInst.inner_transfer}) the inner integrals in
$\mathrm{Dis}(C(u,u',s_c))$ do not change when $h(c)$ is replaced by $g(c)$; since $\nu_M$ is concentrated on
$\mathrm{Good}$, $(C_\infty,\nu_M,g)$ is a disintegration;
\item conditions (i)--(iii) of \cref{pf:sh-local:energyabs} are $(\mathrm{Int})$ and $(\mathrm{Conv})$ at the
diagonalising elements $g(c)$.
\end{itemize}
\Cref{pf:sh-local:energyabs} gives $\int\shlocalHS{O_\infty(g(c))}^2\,d\nu_M(c)\le1/\deg X$. Finally, for
$c\in\mathrm{Good}$, $\shlocalHS{O_\infty(\gamma)}$ is the same for all $\gamma$ diagonalising $c$
(\lref{Formalism.L2.SG.RealObj.InfCut.ooIndep_of_step}, from the same description of diagonalising elements and
the equivariance of $W_u,W_{u'}$ under $S_\infty$); as $\nu_M$ is concentrated on $\mathrm{Good}$, the integral equals
$\shlocalE(u,u',s_c)$.
\end{proof}

\begin{proof}[Proof of \cref{pf:sh-local:realobj}]
By \lref{CompanionTheory.Lime.BM.inputs2_of}, $\sfname{UniqGap}$ gives $(\mathrm{Int})$ and $(\mathrm{Conv})$
for every datum with the gap hypothesis. Let $X,u,u',s_c$ be given. (i) The disintegration identity holds with some
$C>0$ (\lref{Formalism.L2.SG.RealObj.InfCut.ProductExactness.haar_disintegration_real_holds}, the formula for
Haar measure in the coordinates $(s,c,s')$ given by the Gram matrices). (ii) Every $c\in\mathrm{Good}(u,u')$ has a
diagonalising element (\lref{Formalism.L2.SG.RealObj.InfCut.ProductExactness.StepReal_holds}), so $h(c)$
diagonalises $c$, and $(\mathrm{Int})$ gives the integrability at $h(c)$. (iii) is
\cref{pf:sh-local:energystep}.
\end{proof}

% ---------------------------------------------------------------------------
\subsection{Finite dimensionality of spaces of automorphic forms}\label{sec:pf:sh-local:findim}

Let $I$ be an automorphic instance: a totally real number field $F$ with CM extension $E$, a non-degenerate hermitian
matrix $J$ of size $n$ over $E$, an $\OO_E$-lattice $L$, and at each infinite place $v$ an embedding identifying
$\Uni(J)(F_v)$ with $\Uni(J_{\eps_v})$. Let $G_\infty=\prod_v\Uni(J_{\eps_v})$ with Haar measure $\mu_\infty$, $G_f$
the group of finite adelic points of $\Uni(J)$, $K_L\subseteq G_f$ the stabiliser of $L$, $G_\A=G_\infty\times G_f$,
$\Gamma=\Uni(J)(F)\subseteq G_\A$ (embedded diagonally), $\mu_f$ a Haar measure on $G_f$ and
$\mu_\A=\mu_\infty\times\mu_f$. Let $\mathcal F\subseteq G_\A$ be a measurable fundamental domain for $\Gamma$,
$L^2=L^2(\mathcal F,\mu_\A)$, and $R$ the right regular representation of $G_\A$ on $L^2$. For a locally constant
$e:G_f\to\C$ supported in $K_L$ let $R_e$ be the operator with $\langle x,R_ey\rangle=\int e(g)\langle
x,R(1,g)y\rangle\,d\mu_f(g)$.

\begin{theorem}[Finite dimensionality of the row spaces]\label{pf:sh-local:findim}
In this situation, let $\sigma_v$ ($v\mid\infty$) be irreducible unitary representations of $\Uni(J_{\eps_v})$ on
Hilbert spaces, $d_v$ formal degrees for them (Schur orthogonality as in \eqref{eq:pf:sh-local:schur}), and $w_v$ unit
vectors such that $g\mapsto d_v\overline{\langle w_v,\sigma_v(g)w_v\rangle}$ is integrable. Let
$\varphi_\infty(g)=\prod_vd_v\overline{\langle w_v,\sigma_v(g_v)w_v\rangle}$ and $T$ the operator on $L^2$ with
$\langle x,Ty\rangle=\int\varphi_\infty(g)\langle x,R(g,1)y\rangle\,d\mu_\infty(g)$. Then the range of $R_e\circ T$
is finite-dimensional.
\formal{Formalism.L2.SG.RowsDef.RowsFinDim}
\formallabel{CompanionTheory.D.gk.RF.automorphic_forms_finite_dimensional}
\formallabel{CompanionTheory.D.gk.RF.rowsFinDim_of}
\end{theorem}

\begin{remark}
The range of $R_e\circ T$ consists of square-integrable automorphic forms on $\Gamma\backslash G_\A$ of level
$e$ whose components at infinity lie in the $\sigma_v$-isotypic parts generated by the $w_v$. The theorem is a form
of Harish-Chandra's finiteness theorem for spaces of automorphic forms \cite{HarishChandra1968} (compare
\cite{BorelJacquet1979}). It differs from the published statements in two ways: $K$-finiteness and
$\mathfrak z$-finiteness are replaced by the convolution idempotent $\varphi_\infty$ of integrable discrete series
(no $K$-finiteness of $w_v$ is assumed), and the conclusion concerns the closure in $L^2$ of the functions reproduced
by $\varphi_\infty$; the formal proof follows the classical route through reduction theory and cusp estimates.
\end{remark}

\begin{lemma}[Finite dimensionality from linear forms]\label{pf:sh-local:fdlin}
Let $G$ be a second countable locally compact group with Haar measure $\mu$, $\Gamma$ a subgroup, $\mathcal F\subseteq
G$ measurable and $\varphi$ integrable. Call $h:G\to\C$ \emph{reproduced by $\varphi$} if it is bounded, continuous,
left $\Gamma$-invariant and $h(x)=\int\varphi(y)h(xy)\,d\mu(y)$ for all $x$. Assume \emph{detection}: there are a
continuous compactly supported $\alpha$ and $p_1,\dots,p_m\in G$ such that every reproduced $h$ with
$\int_{\mathcal F}|h|^2\,d\mu<\infty$ and $\int\alpha(y)h(p_iy)\,d\mu(y)=0$ for all $i$ vanishes almost everywhere on
$\mathcal F$. Then every subspace of $L^2(\mathcal F,\mu)$ contained in the closure of the classes of reproduced
functions is finite-dimensional.
\formallabel{CompanionTheory.D.gk.RF.finite_dimensional_from_linear_forms}
\end{lemma}

\begin{proof}
Let $\mathcal H_2$ be the vector space of reproduced functions which are square integrable on $\mathcal F$, let
$r:\mathcal H_2\to L^2(\mathcal F)$ be the class map and $l:\mathcal H_2\to\C^m$,
$l(h)=(\int\alpha(y)h(p_iy)\,d\mu(y))_i$ (the integrals converge since $\alpha$ is integrable and $h$ bounded). Detection
says $\ker l\subseteq\ker r$, so $r(\mathcal H_2)$ is a quotient of $\mathcal H_2/\ker l\cong l(\mathcal H_2)\subseteq\C^m$,
hence finite-dimensional and therefore closed. It contains the classes of all reproduced functions which are square
integrable on $\mathcal F$, and a reproduced function whose class lies in $L^2(\mathcal F)$ is such a function; so the
closure of the set of these classes, and every subspace contained in it, lies in $r(\mathcal H_2)$.
\end{proof}

\begin{proof}[Proof of \cref{pf:sh-local:findim}]
The proof establishes seven statements and then derives the theorem from them. Put
$\varphi_0=\varphi_\infty\otimes e$ on $G_\A$.
\begin{enumerate}[label=(\arabic*)]
\item $\varphi_\infty$ is integrable and $\varphi_\infty*\varphi_\infty=\varphi_\infty$, i.e.
$\int\varphi_\infty(y)\varphi_\infty(y^{-1}x)\,d\mu_\infty(y)=\varphi_\infty(x)$
(\lref{Formalism.L2.SG.Pf.phInf_int}, \lref{Formalism.L2.SG.Pf.phInf_conv}; the convolution identity is
Schur orthogonality). The regular representation is unitary with measurable coefficients
(\lref{Formalism.MainNotions.LeviCoordinates.adelic_regular_rep_unitary_continuous}), and $R_e\circ T$ is the
convolution operator by $\varphi_0$ (\lref{Formalism.L2.SG.Pf.prod_op_eq}).
\item \emph{Level.} Since $e$ is locally constant and supported in the compact group $K_L$, there is an open compact
subgroup $K_p\subseteq G_f$ with $e(kx)=e(x)$ for $k\in K_p$ (\lref{CompanionTheory.D.gk.RF.level_structure_of}).
Let $\varphi=\varphi_\infty\otimes\mu_f(K_p)^{-1}\one_{K_p}$; a direct computation with (1) gives
$\varphi*\varphi_0=\varphi_0$.
\item \emph{Cusp estimate} (\lref{CompanionTheory.D.gk.RF.cuspInf} with
\lref{CompanionTheory.D.gk.RF.adelic_cusp_of_infinite_cusp} and
\lref{CompanionTheory.D.gk.RF.congruence_subgroup_compact_open}): for every neighbourhood $U$ of $1$ in $G_\A$
there is a compact $C$ such that every $x\in G_\A$ either lies in $\Gamma C$ or admits a probability measure $\nu$
on $G_\A$ with $x^{-1}mx\in U$ for $\nu$-almost all $m$ and $\int h(mxz)\,d\nu(m)=0$ for all $z\in G_\infty\times
K_p$ and all bounded continuous left $\Gamma$-invariant $h$ \emph{reproduced by $\varphi$}, i.e. with
$h(x)=\int\varphi(y)h(xy)\,d\mu_\A(y)$ for all $x$. (At infinity this is reduction theory for $\Uni(J)$ --- Mahler compactness and
a decomposition along isotropic subspaces --- together with the vanishing of the integrals of $\varphi_\infty$ against
functions invariant under the unipotent radical of a parabolic subgroup.)
\item \emph{Discreteness}: $\Gamma$ meets every compact subset of $G_\A$ in a finite set
(\lref{CompanionTheory.D.gk.RF.fundamental_domain_discrete_of}, \lref{CompanionTheory.D.gk.RF.gamSub}).
\item \emph{Detection} (\lref{CompanionTheory.D.gk.RF.discrete_series_bound_of_inputs}): if $\varphi$ is
integrable, vanishes outside a clopen set $O$ and satisfies the cusp estimate (3) for $O$, and $\Gamma$ is discrete,
then there are a continuous compactly supported $\alpha$ and finitely many $p_1,\dots,p_m\in G_\A$ such that every
bounded continuous $\Gamma$-invariant $h$ reproduced by $\varphi$, with $\int_{\mathcal F}|h|^2<\infty$ and
$\int\alpha(y)h(p_iy)\,d\mu_\A(y)=0$ for all $i$, vanishes in $L^2(\mathcal F)$.
\item \emph{Finite dimensionality}: \cref{pf:sh-local:fdlin}.
\item \emph{The range} (\lref{CompanionTheory.D.gk.RF.convolution_operator_representation},
\lref{CompanionTheory.D.gk.RF.representation_of_convolution}): for $F\in L^2$ essentially bounded on
$\mathcal F$, $R_eTF$ is represented by $x\mapsto\int\varphi_0(y)s(xy)\,d\mu_\A(y)$ with $s$ bounded, measurable and
$\Gamma$-invariant, and such a function is reproduced by $\varphi$ (because $\varphi*\varphi_0=\varphi_0$).
\end{enumerate}
Since essentially bounded functions are dense in $L^2$ (\lref{CompanionTheory.D.gk.RF.dense_bdd}) and
$R_e\circ T$ is continuous, its range lies in the closure of the reproduced functions, and (5), (6) give the claim.
\end{proof}

% ---------------------------------------------------------------------------
\subsection{Isomorphisms \texorpdfstring{$\C\cong\C_q$}{C = Cq}}\label{sec:pf:sh-local:cq}

For a prime $q$ let $\C_q$ be the completion of an algebraic closure of $\Q_q$, and let $\Qbar\subset\C$ be the
field of algebraic numbers in $\C$. We write $\mathfrak c=2^{\aleph_0}$.

\begin{lemma}\label{pf:sh-local:cardCq}
$\C_q$ has cardinality $\mathfrak c$.
\formallabel{GlobalTheory.PadicComplexIso.mk_padicComplex}
\end{lemma}

\begin{proof}
A metric space $Y$ containing a dense subset of cardinality at most $\mathfrak c$ has cardinality at most
$\mathfrak c$, since every point is the limit of a sequence in that subset and
$\mathfrak c^{\aleph_0}=\mathfrak c$ (\lref{GlobalTheory.PadicComplexIso.mk_le_continuum_of_denseRange}).
Since $\Q$ is dense in $\Q_q$, $|\Q_q|\le\mathfrak c$; an algebraic extension of an infinite field $K$ has
cardinality at most $|K|$, so the algebraic closure of $\Q_q$ has cardinality at most $\mathfrak c$; and it is dense in
its completion $\C_q$, so $|\C_q|\le\mathfrak c$. Conversely $\C_q$ is a non-empty complete metric space without
isolated points, hence contains an injective image of $\{0,1\}^{\N}$ (a Cantor set), so $|\C_q|\ge2^{\aleph_0}$.
\end{proof}

\begin{lemma}[Steinitz]\label{pf:sh-local:steinitz}
Let $R$ be a field with $|R|\le\aleph_0$ and $K,L$ algebraically closed fields with $R$-algebra structures such that
$\aleph_0<|K|=|L|$. Then there is a ring isomorphism $e:K\to L$ with $e(a\cdot1_K)=a\cdot1_L$ for all $a\in R$.
\formallabel{GlobalTheory.PadicComplexIso.exists_ringEquiv_over}
\end{lemma}

\begin{proof}
Choose transcendence bases $s$ of $K$ and $t$ of $L$ over $R$. Since $K$ is algebraically closed and
$|R|\le\aleph_0<|K|$, $|K|=|s|$; likewise $|L|=|t|$; so there is a bijection $s\cong t$. It induces an
$R$-isomorphism $R[s]\cong R[t]$ of the subalgebras generated by the bases (both are polynomial rings), and since $K$,
resp. $L$, is an algebraic closure of $R[s]$, resp. $R[t]$, this isomorphism extends to an isomorphism $K\cong L$ by the
uniqueness of algebraic closures. It is $R$-linear by construction.
\end{proof}

\begin{theorem}\label{pf:sh-local:ringequiv}
Let $q$ be a prime and $\lambda:\Qbar\to\C_q$ a ring homomorphism. Then there is a ring isomorphism
$\iota:\C\to\C_q$ with $\iota(x)=\lambda(x)$ for all $x\in\Qbar$.
\formallabel{GlobalTheory.PadicComplexIso.exists_ringEquiv_extending}
\end{theorem}

\begin{proof}
Regard $\C_q$ as a $\Qbar$-algebra via $\lambda$ and $\C$ via the inclusion. $\Qbar$ is algebraic over $\Q$, so
$|\Qbar|\le\max(|\Q|,\aleph_0)=\aleph_0$; $|\C|=\mathfrak c>\aleph_0$; and $|\C_q|=\mathfrak c$ by
\cref{pf:sh-local:cardCq}. \Cref{pf:sh-local:steinitz} with $R=\Qbar$, $K=\C$, $L=\C_q$ gives $\iota$.
\end{proof}

The isomorphism $\iota$ is not continuous and is not unique; the result uses the axiom of choice (through the
existence of transcendence bases), as does the classical statement.

% ---------------------------------------------------------------------------
\subsection{Places of the host field above the small primes}\label{sec:pf:sh-local:host}

\subsubsection{Definitions}
\begin{definition}[Toroidal hosts]\label{pf:sh-local:def:torhost}
Let $E$ be a CM field with ring of integers $\OO_E$ and complex conjugation $a\mapsto\bar a$. A \emph{toroidal host}
$H$ on $E$ consists of a prime $d\equiv7\pmod{16}$, elements $\mathrm{rt},w_0\in\OO_E$ with $\mathrm{rt}^2=-d$ and
$2w_0=1+\mathrm{rt}$, and a complex number $\mathrm{rt}_0$ with $\mathrm{rt}_0^2=-d$; it also records that $\OO_E$ is
unramified over the ring of integers of the maximal real subfield. Its \emph{CM type} is
$\Psi_H=\{\tau:E\to\C:\ \tau(\mathrm{rt})=\mathrm{rt}_0\}$.
\formal{ReciprocityTheory.ChartRestriction.G.Geo.Tor.Host}
\end{definition}

\begin{definition}[Places over $p$ and discriminant depth]\label{pf:sh-local:def:places}
Fix a prime $p$ and put $A_p=\Z_p\otimes\OO_E$, a free $\Z_p$-module of rank $[E:\Q]$, with the involution
$\sigma=1\otimes(\bar{\ })$. A \emph{place of $E$ over $p$} is a primitive idempotent $e$ of $A_p$, i.e.
$e^2=e\ne0$ and every idempotent $f$ with $fe=f$ is $0$ or $e$; it is \emph{split} if $\sigma(e)\ne e$; we write
$\bar e=\sigma(e)$, $A_pe$ for the local factor and $n_e=\rank_{\Z_p}A_pe$. For a commutative $\Z_p$-algebra $O$ and
a ring homomorphism $\tau:\OO_E\to O$ let $\tilde\tau:A_p\to O$ be the $\Z_p$-algebra map with
$\tilde\tau(1\otimes a)=\tau(a)$, and $\mathrm{Emb}(e)=\{\tau:\tilde\tau(e)=1\}$ the set of embeddings
\emph{belonging to} $e$.

In the following, $C$ is an algebraically closed field of characteristic $0$ with a valuation $v:C\to\R_{\ge0}$
(multiplicative notation) whose ring of integers is a given ring $O$ (with an injective map $O\to C$ onto
$\{v\le1\}$), $p$ is not a unit in $O$, and $O$ carries a $\Z_p$-algebra structure. For a split place $e$, a bijection
$t:\{1,\dots,n_e\}\to\mathrm{Emb}(e)$ and $v_d\in\R$, we say that $(e,t)$ has \emph{discriminant depth} $v_d$ if
there is a $\Z_p$-basis $(b_i)$ of $A_pe$ with
\[
v\bigl(\det(\tilde t_j(b_i))_{i,j}^2\bigr)=v(p)^{n_ev_d}.
\]
\formal{ReciprocityTheory.ChartRestriction.G.Pdiv.AvDefs.Place}
\formal{ReciprocityTheory.ChartRestriction.G.Pdiv.AvDefs.PlaceU}
\formal{ReciprocityTheory.ChartRestriction.G.Pdiv.AvDefs.Emb}
\formal{ReciprocityTheory.ChartRestriction.G.Pdiv.AvDefs.HasVd}
\end{definition}

\begin{definition}[Split places]\label{pf:sh-local:def:splitplaces}
Let $H$ be a toroidal host on $E$, $p$ a prime, $f\in\N$ and $v_d\in\R$. Then $d$ \emph{splits at $p$},
$\mathrm{SplitAt}(d,p)$, if $p=2$ or $\bigl(\frac{-d}p\bigr)=1$. The \emph{split-places condition}
$\mathrm{SplitPlaces}(p,H,f,v_d)$ says: for all $C,v,O$ as above and every injective map $\tau\mapsto\tau_O$ from
the embeddings $E\to\C$ to the ring homomorphisms $\OO_E\to O$ such that $\tau(\mathrm{rt})=\tau'(\mathrm{rt})$ if
and only if $\tau_O(\mathrm{rt})=\tau'_O(\mathrm{rt})$, there are split places $e_1,\dots,e_s$ of $E$ over $p$ and
real numbers $v_{d,1},\dots,v_{d,s}$ such that
\begin{enumerate}[label=(\roman*)]
\item $\sum_in_{e_i}=f$ and $\sum_in_{e_i}v_{d,i}=f\,v_d$;
\item every element of $\mathrm{Emb}(e_i)$ is $\tau_O$ for some $\tau\in\Psi_H$, and every element of
$\mathrm{Emb}(\bar e_i)$ is $\tau_O$ for some $\tau\notin\Psi_H$;
\item for every $\tau\in\Psi_H$ there is exactly one $i$ with $\tau_O\in\mathrm{Emb}(e_i)$;
\item for every $i$ there is a bijection $t_i:\{1,\dots,n_{e_i}\}\to\mathrm{Emb}(e_i)$ such that $(e_i,t_i)$ has
discriminant depth $v_{d,i}$.
\end{enumerate}
\formal{UpperTheory.Depth.SplitAt}\formal{UpperTheory.Depth.SplitPlaces}
\end{definition}

\subsubsection{Three lemmas on idempotents}

\begin{lemma}[The split idempotent]\label{pf:sh-local:splitidem}
Let $H$ be a toroidal host on $E$ and $p$ a prime with $\mathrm{SplitAt}(d,p)$. There are an idempotent $e\in A_p$
and $\alpha\in\Z_p$ with $e\ne0,1$, $e+\sigma(e)=1$ and $\alpha^2=-d$, such that for every domain $O$ of
characteristic $0$ with a $\Z_p$-algebra structure and every ring homomorphism $\tau:\OO_E\to O$:
$\tau\in\mathrm{Emb}(e)\iff\tau(\mathrm{rt})=\alpha$, and $\tau\in\mathrm{Emb}(1-e)\iff\tau(\mathrm{rt})=-\alpha$.
\formallabel{UpperTheory.exists_splitIdempotent_of_splitAt}
\end{lemma}

\begin{proof}
Write $d+1=8k$ (possible as $d\equiv7\pmod{16}$). \emph{Step 1: there are $\beta,\gamma\in\Z_p$ with
$\beta^2=\beta-2k$ and $\gamma(1-2\beta)=1$} (\lref{UpperTheory.Depth.SplitIdem.exists_root}). If $p=2$,
Hensel's lemma applied to $X^2-X+2k$ at $X=0$ (the value $2k$ has norm $<1$ and the derivative $-1$ is a unit)
gives $\beta$, and $1-2\beta$ is a unit of $\Z_2$. If $p\ne2$, then $\bigl(\frac{-d}p\bigr)=1$ gives $a\in\Z$ with
$a^2\equiv-d\pmod p$ and $p\nmid a$; Hensel's lemma applied to $X^2+d$ at $a$ (the derivative $2a$ is a unit) gives
$\alpha_0\in\Z_p$ with $\alpha_0^2=-d$. With $u=2^{-1}$ and $v_0=(-d)^{-1}$ in $\Z_p$ put $\beta=u(1-\alpha_0)$ and
$\gamma=v_0\alpha_0$; then $\beta^2-\beta=(\alpha_0^2-1)/4=-2k$ and $\gamma(1-2\beta)=v_0\alpha_0^2=1$.

\emph{Step 2.} In $\OO_E$, $w_0^2=w_0-2k$ (expand $(2w_0-1)^2=\mathrm{rt}^2=1-8k$), and $\bar w_0=1-w_0$: indeed
$\overline{\mathrm{rt}}=\pm\mathrm{rt}$, and $\overline{\mathrm{rt}}=\mathrm{rt}$ is impossible because then
$\tau(\mathrm{rt})$ would be real for every embedding $\tau$, whereas $\tau(\mathrm{rt})^2=-d<0$. Let
$e=\gamma(w_0-\beta)\in A_p$ (identifying $\Z_p$ and $\OO_E$ with their images). Then
$e+\sigma(e)=\gamma(w_0-\beta+1-w_0-\beta)=\gamma(1-2\beta)=1$, and from
$(w_0-\beta)^2=w_0-2k-2\beta w_0+\beta-2k=(w_0-\beta)(1-2\beta)$ we get $e^2=\gamma(w_0-\beta)=e$. As $A_p\ne0$
($\Z_p$ is flat over $\Z$), $e+\sigma(e)=1$ excludes $e=0$ and $e=1$. Put $\alpha=1-2\beta$; then
$\alpha^2=1-4\beta+4\beta^2=1-8k=-d$.

\emph{Step 3.} Let $\tau:\OO_E\to O$. Then $\tilde\tau(e)=\gamma(\tau(w_0)-\beta)$ and $2\tau(w_0)=1+\tau(\mathrm{rt})$.
Since $\gamma(1-2\beta)=1$ and $2$ is not a zero divisor in $O$, $\tilde\tau(e)=1$ if and only if
$\tau(w_0)=1-\beta$, i.e. $\tau(\mathrm{rt})=1-2\beta=\alpha$; and $\tilde\tau(1-e)=1$ if and only if
$\tau(w_0)=\beta$, i.e. $\tau(\mathrm{rt})=-\alpha$.
\end{proof}

\begin{lemma}[Primitive decomposition]\label{pf:sh-local:primdec}
In a noetherian commutative ring, every idempotent is a finite sum of pairwise orthogonal primitive idempotents.
\formallabel{UpperTheory.exists_primitive_orthogonal_idempotents_sum_eq_of_noetherian}
\end{lemma}

\begin{proof}
We use noetherian induction on the ideal $(1-e)$. If $e=0$ take the empty sum, and if $e$ is primitive take $e$
itself. Otherwise there is an idempotent $x$ with $xe=x$ and $x\ne0,e$; then $x$ and $e-x$ are orthogonal
idempotents with sum $e$, and the ideals $(1-x)$ and $(1-(e-x))$ strictly contain $(1-e)$
(\lref{idem_lt}). By induction both are sums of orthogonal primitive idempotents; the elements of
the two decompositions are mutually orthogonal (they are absorbed by $x$, resp. $e-x$, and $x(e-x)=0$), so the
union is a decomposition of $e$.
\end{proof}

\begin{lemma}[Counting embeddings]\label{pf:sh-local:countemb}
Let $C,v,O$ be as above (with any $\Z_p$-algebra structure on $O$) and $e\in A_p$ an idempotent. Then there is a
bijection $\{1,\dots,\rank_{\Z_p}A_pe\}\cong\mathrm{Emb}(e)$, and there are exactly $[E:\Q]$ ring homomorphisms
$\OO_E\to O$.
\formallabel{UpperTheory.trace_form_decomposition_CMDatum}
\end{lemma}

\begin{proof}
Every field embedding $E\to C$ maps $\OO_E$ into the integrally closed ring $O$, and every ring homomorphism
$\OO_E\to O$ extends uniquely to $E$; so the ring homomorphisms $\OO_E\to O$ correspond to the $[E:\Q]$ embeddings
$E\to C$ (\lref{UpperTheory.TraceSumEquivalence.exists_equiv}). Multiplication by $e$ is an idempotent endomorphism of the free
$\Z_p$-module $A_p$, so its trace is the rank of its image $A_pe$. On the other hand the trace of multiplication by
an element $x$ of $A_p$ equals $\sum_\tau\tilde\tau(x)$, the sum over all ring homomorphisms $\tau:\OO_E\to O$
(\lref{UpperTheory.TraceSumEquivalence.trace_eq_sum}). Each $\tilde\tau(e)$ is an idempotent of the domain $O$, hence $0$ or $1$, so the
trace is $|\mathrm{Emb}(e)|$.
\end{proof}

\subsubsection{The discriminant as a product over places}

\begin{proposition}\label{pf:sh-local:discprod}
Let $C,v,O$ be as above and $e_1,\dots,e_s$ split places of $E$ over $p$ with $e_ie_j=0$ for $i\ne j$,
$e_i\bar e_j=0$ for all $i,j$, and $\sum_ie_i+\sum_i\bar e_i=1$. Then $2\sum_in_{e_i}=[E:\Q]$, and for all bijections
$t_i:\{1,\dots,n_{e_i}\}\to\mathrm{Emb}(e_i)$ there are $v_{d,i}\in\R$ such that $(e_i,t_i)$ has discriminant depth
$v_{d,i}$ for every $i$ and
\[
\sum_in_{e_i}v_{d,i}=\tfrac12\,v_p(\mathrm{disc}\,E).
\]
\formallabel{UpperTheory.discr_CMDatum_eq_prod_places}
\end{proposition}

\begin{proof}
\emph{Ranks.} $A_p$ is the direct sum of the $A_pe_i$ and $A_p\bar e_i$
(\lref{UpperTheory.PlacesDiscr.splitEquiv}), and $\sigma$ maps $A_pe_i$ isomorphically onto $A_p\bar e_i$;
comparing ranks, $[E:\Q]=\rank_{\Z_p}A_p=2\sum_in_{e_i}$.

\emph{The block matrix.} Choose $\Z_p$-bases $B_i=(B_{i,a})_a$ of $A_pe_i$ and let
$M_i=(\tilde t_i(b)(B_{i,a}))_{a,b}\in M_{n_{e_i}}(O)$. The ring homomorphisms $\OO_E\to O$ are exactly the
$t_i(b)$ and the $t_i(b)\circ(\bar{\ })$, each once: for every $\tau$ the elements $\tilde\tau(e_i)$,
$\tilde\tau(\bar e_i)$ are idempotents of the domain $O$ summing to $1$ and pairwise orthogonal, so exactly one of
them is $1$, and $\widetilde{\tau\circ\bar{\ }}(e_i)=\tilde\tau(\bar e_i)$; with the count of
\cref{pf:sh-local:countemb} this is a bijection. Order the basis $\gamma$ of $A_p$ formed by the $B_i$ and the
$\sigma(B_i)$, and the homomorphisms, compatibly. If $\tau\in\mathrm{Emb}(e_i)$ then $\tilde\tau$ vanishes on $A_pe_j$
for $j\ne i$ and on all $A_p\bar e_j$, because $e_ie_j=0=e_i\bar e_j$; hence the matrix
$(\tilde\tau(\gamma_m))_{m,\tau}$ is block diagonal with blocks $\mathrm{diag}(M_i,M_i)$ (for the second block use
$\widetilde{\tau\circ\bar{\ }}(\sigma(x))=\tilde\tau(x)$), and its determinant is $\prod_i\det(M_i)^2$.

\emph{Comparison with the discriminant.} Let $(\omega_m)$ be an integral basis of $\OO_E$; then $(1\otimes\omega_m)$
is another $\Z_p$-basis of $A_p$, and the change of basis matrix to $\gamma$ has a unit determinant in $\Z_p$, whose
image in $O$ has valuation $1$. Since $\mathrm{disc}\,E=\det(\tau(\omega_m))^2$ (the embeddings $E\to C$ being
identified with the homomorphisms $\OO_E\to O$), we obtain
$\prod_iv(\det(M_i)^2)^2=v(\mathrm{disc}\,E)$. For a non-zero integer $N$ one has $v(N)=v(p)^{v_p(N)}$, since the
primes $\ne p$ are units of $\Z_p$ (\lref{valuation_intCast_eq_pow_padicValInt}).

\emph{Depths.} Put $Q=v(p)$, so $0<Q<1$, and $x_i=v(\det(M_i)^2)>0$. Let $v_{d,i}=\log_Q(x_i)/n_{e_i}$ (and $0$ if
$n_{e_i}=0$, in which case $x_i=1$). Then $x_i=Q^{n_{e_i}v_{d,i}}$, so $(e_i,t_i)$ has discriminant depth $v_{d,i}$
with the basis $B_i$, and $\sum_in_{e_i}v_{d,i}=\log_Q\prod_ix_i=\frac12\log_Q\prod_ix_i^2=\frac12v_p(\mathrm{disc}\,E)$.
\end{proof}

\begin{proposition}\label{pf:sh-local:splitplaces}
Let $H$ be a toroidal host on $E$, $p$ a prime with $\mathrm{SplitAt}(d,p)$ and $f\in\N$ with $[E:\Q]=2f$. Then
$\mathrm{SplitPlaces}\bigl(p,H,f,v_p(\mathrm{disc}\,E)/[E:\Q]\bigr)$ holds.
\formallabel{UpperTheory.splitPlaces_discriminantDepth_of_splitAt}
\end{proposition}

\begin{proof}
Let $C,v,O$ and $\tau\mapsto\tau_O$ be given. $O$ is a domain of characteristic $0$. There is $\tau_0\in\Psi_H$
(\lref{UpperTheory.Depth.SplitPl.exists_inPsi}). By \cref{pf:sh-local:countemb} both sets have $[E:\Q]$
elements, so the injective map $\tau\mapsto\tau_O$ is bijective. By the compatibility hypothesis,
$\tau\in\Psi_H$ if and only if $\tau_O(\mathrm{rt})=\tau_{0,O}(\mathrm{rt})$; and if
$\tau_O(\mathrm{rt})=-\tau_{0,O}(\mathrm{rt})$ then $\tau\notin\Psi_H$, for otherwise
$2\tau_{0,O}(\mathrm{rt})=0$, whereas $\tau_{0,O}(\mathrm{rt})^2=-d\ne0$ in $O$.

Take $e,\alpha$ from \cref{pf:sh-local:splitidem}. Since $\alpha^2=-d=\tau_{0,O}(\mathrm{rt})^2$, the image of $\alpha$ in
$O$ is $\pm\tau_{0,O}(\mathrm{rt})$; replacing $(e,\alpha)$ by $(1-e,-\alpha)$ if necessary we get an idempotent $e$ with
$\sigma(e)=1-e$, $\mathrm{Emb}(e)=\{\tau':\tau'(\mathrm{rt})=\tau_{0,O}(\mathrm{rt})\}$ and
$\mathrm{Emb}(1-e)=\{\tau':\tau'(\mathrm{rt})=-\tau_{0,O}(\mathrm{rt})\}$. As $A_p$ is noetherian,
\cref{pf:sh-local:primdec} writes $e=\sum_{i\le s}g_i$ with pairwise orthogonal primitive idempotents $g_i$. Then
$g_ie=g_i$, $\sigma(g_i)(1-e)=\sigma(g_i)$, hence $g_i\sigma(g_j)=g_ie\,\sigma(g_j)(1-e)=0$ for all $i,j$; in
particular $\sigma(g_i)\ne g_i$, so the $g_i$ are split places, and $\sum g_i+\sum\sigma(g_i)=e+\sigma(e)=1$.

\Cref{pf:sh-local:discprod} gives $2\sum_in_{g_i}=[E:\Q]=2f$, and, with bijections
$t_i:\{1,\dots,n_{g_i}\}\to\mathrm{Emb}(g_i)$ from \cref{pf:sh-local:countemb}, numbers $v_{d,i}$ with (iv) and
$\sum_in_{g_i}v_{d,i}=\frac12v_p(\mathrm{disc}\,E)=f\cdot v_p(\mathrm{disc}\,E)/[E:\Q]$; this is (i). For (ii): if
$xy=x$ and $\tilde\tau(x)=1$ then $\tilde\tau(y)=1$; so $\tau'\in\mathrm{Emb}(g_i)$ implies $\tau'\in\mathrm{Emb}(e)$,
i.e. $\tau'(\mathrm{rt})=\tau_{0,O}(\mathrm{rt})$, and writing $\tau'=\tau_O$ we get $\tau\in\Psi_H$; similarly
$\tau'\in\mathrm{Emb}(\sigma(g_i))$ implies $\tau'\in\mathrm{Emb}(1-e)$, hence $\tau'=\tau_O$ with $\tau\notin\Psi_H$.
For (iii): if $\tau\in\Psi_H$ then $\tilde\tau_O(e)=1$, and the values $\tilde\tau_O(g_i)$ are orthogonal idempotents of
the domain $O$ with sum $1$, so exactly one of them is $1$ (\lref{UpperTheory.Depth.SplitPl.existsUnique_one}).
\end{proof}

\subsubsection{The host of a datum}
Let $D$ be a datum with host (\cref{def:host}) consisting of primes $q,d$, the integer $Q_S$ and the CM field $E$ of
degree $2f$; let $S_0$ be the set of prime divisors of $Q_S$ (the \emph{small primes}),
$v_{dH}(p)=v_p(\mathrm{disc}\,E)/[E:\Q]$ and $v_{d2}=\log_2Q_S$.

\begin{theorem}[Places of the host over the small primes]\label{pf:sh-local:twoplaces}
For every datum $D$, every toroidal host $H$ on the host field $E$ of $D$ whose prime $d$ is the prime $d$ of the
host of $D$, and every prime $p\in S_0$: $\mathrm{SplitAt}(d,p)$, $\mathrm{SplitPlaces}(p,H,f,v_{dH}(p))$, and
$0\le v_{dH}(p)\le v_{d2}$.
\formallabel{GlobalTheory.CharDepth.twoPlaces_host}
\end{theorem}

\begin{proof}
If $p=2$, $d$ splits at $p$ by definition; if $p\ne2$, then $p\mid Q_S$ and the host condition
$\bigl(\frac{-d}l\bigr)=1$ for odd $l\mid Q_S$ (\cref{def:host}) gives $\mathrm{SplitAt}(d,p)$. By
\lref{UpperTheory.cmDatum_discriminantDepth_bounds}, for every prime $p$ the number $v_{dH}(p)$ equals
$v_p(\mathrm{disc}\,E)/[E:\Q]$ (with the formal definition through $|\mathrm{disc}\,E|$), and for $p\in S_0$ one has
$v_{dH}(p)\le v_p(Q_S)\le\log_2Q_S$; clearly $v_{dH}(p)\ge0$. Since $[E:\Q]=2f$
(\lref{GlobalTheory.PlacesOverG.finrank_E'}), \cref{pf:sh-local:splitplaces} gives
$\mathrm{SplitPlaces}(p,H,f,v_{dH}(p))$.
\end{proof}

% ---------------------------------------------------------------------------
\subsection{Characters of formal groups with jet charts}\label{sec:pf:sh-local:chart}

This subsection proves the three statements $(\mathrm H_p)$, $(\mathrm A_p)$, $(\mathrm U_p)$ used in
\cref{sec:pf:onto1,sec:pf:onto2} to bound depths of theta exponents at the primes over $p$. All three are
obtained from a common linear-algebraic model of the formal group of a degenerating abelian variety (the
\emph{linear data} below), and then from three statements about linear data.

\subsubsection{The setting}\label{sec:pf:sh-local:chartsetting}
\begin{definition}[Jet charts and characters]\label{pf:sh-local:def:jetchart}
\emph{The base.} Let $p$ be a prime and $T$ a complete noetherian local domain of characteristic $0$ in which $p$
is not a unit. Let $G$ be a commutative group scheme over $T$, $\mathcal D$ a degeneration datum for the maximal
ideal of $T$ (with abelian part an abelian scheme $P_{\mathcal D}$ over $T$ of relative dimension $g$, and a torus
part with character group $X$), $A$ a polarised abelian variety over the fraction field of $T$, and $U$ a
uniformisation datum relating $G$, $\mathcal D$ and $A$ (in the sense of Mumford's construction
\cite[Ch.~III]{FaltingsChai1990}).

\emph{Jet charts.} Let $d\in\N$ and $F\in T[[X,Y]]^d$ a formal group law of dimension $d$. For $m\in\N$ let $J_m$ be
the ring of $m$-jets $T[X_1,\dots,X_d]/(X_1,\dots,X_d)^{m+1}$. A \emph{jet chart} $x$ of $G$ for $F$ is a family of
morphisms $x_m:\Spec J_m\to G$ over $T$ compatible with truncation, sending the zero jet to the unit section,
universal among jets at the unit section (every morphism from the spectrum of an $m$-jet ring $T[Y_\kappa]/(Y)^{m+1}$
to $G$ that sends the zero jet to the unit factors uniquely through $x_m$ by a jet homomorphism), and transforming
$F$ into the group law: $x_m\circ F^*=x_m\circ\mathrm{pr}_0\cdot x_m\circ\mathrm{pr}_1$.

\emph{Coefficients.} Let $C$ be an algebraically closed field of characteristic $0$ with a valuation $v$ (values in
$\R_{\ge0}$), $O$ its ring of integers (via a given injection $O\to C$), such that $p$ is not a unit of $O$ and $O$ is
$p$-adically complete, and $\iota:T\to O$ a local homomorphism. Let $F^\iota$ be the image of $F$ under $\iota$. A
\emph{character} of $F^\iota$ is a power series $u\in O[[X_1,\dots,X_d]]$ with $u(0)=1$ and
$u(F^\iota(X,Y))=u(X)u(Y)$; its \emph{logarithmic differential} $du\in O^d$ is the vector of its coefficients of
$X_1,\dots,X_d$.
\formal{ReciprocityTheory.ChartRestriction.G.Out.UnifPre}
\formal{ReciprocityTheory.ChartRestriction.G.Pdiv.ChartF.IsChartJ}
\formal{ReciprocityTheory.ChartRestriction.G.Pdiv.ChartF.IsCh}
\end{definition}

\begin{definition}[Actions, eigenbases and valuation lengths]\label{pf:sh-local:def:eigen}
\emph{Actions.} For $(\mathrm A_p)$ and $(\mathrm U_p)$ we are moreover given a commutative ring $\OO_E$ with
involution $*$ which is a Dedekind domain of characteristic $0$, free of finite rank over $\Z$; an $\OO_E$-module
structure on $X$; actions of $\OO_E$ on $G$, on $\mathcal D$ and on $A$ for which $U$ is $\OO_E$-linear; and an action
$a\mapsto\alpha(a)$ of $\OO_E$ on $F$ by $d$-tuples of power series without constant term, compatible with the chart:
$x_m\circ\alpha(a)^*$ is the action of $a$ on $G$ composed with $x_m$, for all $m$. Let $\rho(a)\in M_d(O)$ be the
linear part of $\alpha(a)^\iota$.

\emph{Eigenbases and lengths.} For $\tau:\OO_E\to O$, an \emph{eigenbasis} of rank $s$ at $\tau$ is a family
$W=(W_1,\dots,W_s)$ in $O^d$ with $\rho(a)W_i=\tau(a)W_i$ for all $a$, such that every $y\in O^d$ with
$\rho(a)y=\tau(a)y$ for all $a$ is a unique $O$-linear combination of the $W_i$. Let
$\mathrm{res}_W:O^d\to O^s$, $f\mapsto(\sum_jW_{i,j}f_j)_i$. An $O$-module $M$ has \emph{valuation length}
$x\in\R_{\ge0}$ if $M\cong\bigoplus_{i=1}^mO/(c_i)$ with $\prod_iv(c_i)=x$ (a multiplicative version of the length).
Finally $\mathrm{AF}(W,x)$ means: there are characters $u_1,\dots,u_N$ of $F^\iota$ such that
$O^s/\langle\mathrm{res}_W(du_k)\rangle_O$ has valuation length $x$.
\formal{ReciprocityTheory.ChartRestriction.G.Pdiv.ChartF.ChartActs}
\formal{ReciprocityTheory.PGsh.Up.HasLen}
\formal{ReciprocityTheory.ChartRestriction.G.Pdiv.ChartF.HasAF}
\end{definition}

The linear part $\rho$ is a ring homomorphism $\OO_E\to M_d(O)$
(\lref{GlobalTheory.T.formalGroupAction_linearPart_ringHom}).

\begin{definition}[The three statements]\label{pf:sh-local:def:three}
Let $p$ be a prime.
\begin{enumerate}[label=(\alph*)]
\item $(\mathrm H_p)$: for all data of the base, jet charts and coefficients as above, there are characters
$u_1,\dots,u_N$ of $F^\iota$ such that $c\,f\in\langle du_1,\dots,du_N\rangle_O$ for all $f\in O^d$ and all $c\in O$
with $v(c)^{p-1}\le v(p)$.
\item $(\mathrm A_p)$: for all data as above with an action of $\OO_E$, every $\Z_p$-algebra structure on $O$, every
split place $w$ of $E$ over $p$ (\cref{sec:pf:sh-local:host}, with $n=n_w$), every bijection
$t:\{1,\dots,n\}\to\mathrm{Emb}(w)$ such that $(w,t)$ has discriminant depth $v_d$, all $r_p,r_m\in\N$, all eigenbases
$W_j$ of rank $r_p$ at $t_j$ ($1\le j\le n$) such that every $\tau\in\mathrm{Emb}(\bar w)$ has an eigenbasis of rank
$r_m$: there are $x_1,\dots,x_n$ with $\mathrm{AF}(W_j,x_j)$ for all $j$ and
\[
v(p)^{\,n\,r_m\left(\frac1{p-1}+\frac{v_d}2\right)}\le\prod_jx_j .
\]
\item $(\mathrm U_p)$: assume $p\ne2$. For all data as above with an action of $\OO_E$ such that $p$ does not divide
the discriminant of $\OO_E$, every $\Z_p$-algebra structure on $O$ ($O$ local), every place $w$ of $E$ over $p$ with
$n=n_w\ge1$, every bijection $t:\Z/n\to\mathrm{Emb}(w)$ such that some $\Z_p$-basis $(b_i)$ of $A_pw$ has
$\det(\tilde t_j(b_i))$ a unit and $t_{j+1}(a)\equiv t_j(a)^p$ modulo the maximal ideal of $O$ for all $j,a$, every
$r:\Hom(\OO_E,O)\to\N$ and eigenbases $W(\tau)$ of rank $r(\tau)$ at every $\tau$, and every $j$: there is $x$ with
$\mathrm{AF}(W(t_j),x)$ and
\[
v(p)^{\sum_{l\ge0}p^{-(l+1)}\,r(t_{j+l+1}\circ*)}\le x .
\]
\end{enumerate}
\formal{ReciprocityTheory.ChartRestriction.G.Pdiv.ChartF.HTFSt}
\formal{ReciprocityTheory.ChartRestriction.G.Pdiv.ChartF.AvFSt}
\formal{ReciprocityTheory.ChartRestriction.G.Pdiv.ChartF.UpFSt}
\end{definition}

\begin{theorem}\label{pf:sh-local:Hp}
$(\mathrm H_p)$ holds for every prime $p$.
\formallabel{GlobalTheory.ChF.hodgeTate_inclusion_chart}
\end{theorem}

\begin{theorem}\label{pf:sh-local:Up}
$(\mathrm U_p)$ holds for every prime $p$.
\formallabel{GlobalTheory.ChF.unramified_companion_chart_characters}
\end{theorem}

\begin{theorem}\label{pf:sh-local:Ap}
$(\mathrm A_p)$ holds for every prime $p$.
\formallabel{GlobalTheory.ChF.colength_bound_split_place_chart}
\end{theorem}

\subsubsection{Linear data}
Let $\OO_E$ be a commutative ring with involution $*$, $\sigma=1\otimes*$ on $\Z_p\otimes\OO_E$, and let $c_p\in O$
and $H$ be a Hodge--Tate setting and datum of genus $g$ for $(O,\Z_p\otimes\OO_E,\sigma,c_p)$ (\cref{pf:tower:def-htdatum};
we write $\rho_H,\rho_{L},P,P_L,E,H_t,B_t$ for its components).

\begin{definition}[Linear data]\label{pf:sh-local:def:lindat}
Let $d,r,m\in\N$. \emph{Linear data} for $H$ consist of matrices $A\in M_{d\times r}(O)$, $B\in M_{g\times d}(O)$,
ring homomorphisms $\rho_G:\OO_E\to M_d(O)$ and $N:\OO_E\to M_r(\Z)$, and vectors $s_1,\dots,s_m\in O^d$, such that:
$P=0$ and $P_L=0$; the sequence $O^r\xrightarrow{A}O^d\xrightarrow{B}O^g$ is split exact, i.e. $BA=0$ and there are
$A',B'$ with $A'A=1$, $BB'=1$, $AA'+B'B=1$; $B\rho_G(a)=\rho_L(1\otimes a)B$ and $\rho_G(a)A=AN(a)$ for all $a$;
for every index $i$ of the columns of $H_t$ some $s_k$ equals $B^{\mathsf T}(H_t)_{\cdot,i}$; and for every
$i\le r$ some $s_k$ satisfies $A^{\mathsf T}s_k=\delta_i$ (the $i$-th standard vector).
\formal{GlobalTheory.ChF.LinDat}
\end{definition}

In words: $O^d$ is the Lie algebra of the formal group, $A$ embeds the Lie algebra of the toric part, $B$ projects to
the Lie algebra of the abelian part, $H_t$ is the Hodge--Tate map of the abelian part, and the $s_k$ are logarithmic
differentials of characters which restrict to a basis of characters of the torus and pull back the Hodge--Tate
characters of the abelian part.

\begin{proposition}[Linear data exist]\label{pf:sh-local:lindatex}
Let the base, jet chart and coefficients be as in \cref{sec:pf:sh-local:chartsetting}, with an action of a
commutative ring $\OO_E$ with involution as there (for $(\mathrm H_p)$, $\OO_E=\Z$ acting by multiplication), and fix
a $\Z_p$-algebra structure on $O$. Then there are $c_p$, a Hodge--Tate setting and a Hodge--Tate datum $H$ of genus
$g$ for $(O,\Z_p\otimes\OO_E,\sigma,c_p)$, $r,m\in\N$, linear data $(A,B,\rho,N,(du_k)_{k\le m})$ for $H$ with
$\rho_G=\rho$ the linear part of the action on $F^\iota$, where $u_1,\dots,u_m$ are characters of $F^\iota$, and
$A$ is the linear part of a $d$-tuple of power series in $r$ variables.
\end{proposition}

\begin{proof}
This is the common part of the three formal proofs; all steps are results of the development.
\begin{enumerate}[label=(\arabic*)]
\item \emph{Nilpotent charts.} The jet chart $x$ gives a chart $X$ of $G$ on nilpotent coordinates for $F$, under
which the action of each $a\in\OO_E$ on $G$ is read by the power series $\alpha(a)$
(\lref{GlobalTheory.T.exists_nilChart_reads_of_isChartJ}). The abelian scheme $P_{\mathcal D}$, smooth of
relative dimension $g$, has a formal group law $F_B$ with a nilpotent chart $X_B$
(\lref{GlobalTheory.T.exists_nilpotentGroupChart_smooth_commGroup}), and every endomorphism, in particular the
action of $a$ on $P_{\mathcal D}$, is read by power series $\alpha_B(a)$
(\lref{GlobalTheory.T.exists_powerSeries_reading_groupHom_nilpotentCharts}).
\item \emph{The uniformisation splits $F$.} There are $r\in\N$, $J\in T[[Y_1,\dots,Y_r]]^d$ and
$Q\in T[[X_1,\dots,X_d]]^g$ exhibiting $F$ as an extension of $F_B$ by the formal torus of dimension $r$, and a ring
homomorphism $N:\OO_E\to M_r(\Z)$ with $\mathrm{lin}(Q)\,\mathrm{lin}(\alpha(a))=\mathrm{lin}(\alpha_B(a))\,\mathrm{lin}(Q)$ and
$\mathrm{lin}(\alpha(a))\,\mathrm{lin}(J)=\mathrm{lin}(J)\,N(a)$, $\mathrm{lin}$ denoting the matrix of linear parts
(\lref{GlobalTheory.T.uniformization_split_formal_group_law_equivariant}).
\item \emph{Change of coefficients.} Base change along $\iota$ preserves formal group laws, linear parts and the
extension property (\lref{GlobalTheory.T.formalGroupLaw_jacobian_extension_changeCoeff}).
\item \emph{Torsion.} For every $m$ the $p^m$-torsion algebra $O[[X]]/([p^m]_{F^\iota})$ is finite free over $O$,
reduced after $\otimes_OC$, and its $O$-points separate its elements
(\lref{GlobalTheory.FormalTorsion.fhSt}).
\item \emph{Hodge--Tate datum of the abelian part.} There are $c_p$, a Hodge--Tate setting and datum $H$ of genus
$g$ with $P=P_L=0$, an invertible $S\in M_g(O)$ with $\mathrm{lin}(\alpha_B(a)^\iota)S=S\rho_L(1\otimes a)$, and for every
column index $i$ a character $u$ of $F_B^\iota$ with $S^{\mathsf T}du=(H_t)_{\cdot,i}$
(\lref{GlobalTheory.T.exists_hodgeTateDatum_polarisedAbelian_characters}; compare \cref{sec:pf:tower}).
\item \emph{Linear data.} From (2)--(5), \lref{GlobalTheory.T.exists_linearData_of_split_formalGroup_extension}
gives $c_p,H$ as in (5), a matrix $B$, $m$ and characters $u_1,\dots,u_m$ of $F^\iota$ such that
$(\mathrm{lin}(J^\iota),B,\rho,N,(du_k))$ are linear data for $H$.
\end{enumerate}
For $(\mathrm H_p)$ the $\Z_p$-algebra structure on $O$ exists because $O$ is $p$-adically complete
(\lref{GlobalTheory.ChF.padicInt_algebra_of_isAdicComplete_and_matrix_unique}, which also shows that every ring
homomorphism $\Z_p\to M_n(O)$ is the scalar one), and the action of $\Z$ on the polarised abelian part by
multiplication is compatible with the Rosati involution by the theorem of the cube
(\lref{GlobalTheory.exists_oAct_int}, \lref{GlobalTheory.firstFiveLaws}); the intertwining relations of
(2) for $a\in\Z$ are the identities $\mathrm{lin}(\alpha(a))=a\cdot1$.
\end{proof}

\subsubsection{The Hodge--Tate inclusion}

\begin{lemma}\label{pf:sh-local:htsurj}
Let $R$ be a domain, $E\in M_{2g}(R)$ alternating with unit determinant, $c_p\in R\setminus\{0\}$ and
$H_t,B_t\in M_{g\times2g}(R)$ with $B_t^{\mathsf T}H_t-H_t^{\mathsf T}B_t=c_pE$. Then for every $f\in R^g$ there is
$t\in R^{2g}$ with $H_tt=c_pf$.
\formallabel{ReciprocityTheory.PGsh.Up.ht_cp_surj}
\end{lemma}

\begin{proof}
Let $\Phi=\binom{H_t}{B_t}\in M_{2g}(R)$ and $\Psi=E^{-1}(B_t^{\mathsf T},\,-H_t^{\mathsf T})$. Then
$\Psi\Phi=E^{-1}(B_t^{\mathsf T}H_t-H_t^{\mathsf T}B_t)=c_p\cdot1$. Over the fraction field this shows that $\Phi$ is
invertible with $\Psi=c_p\Phi^{-1}$, hence also $\Phi\Psi=c_p\cdot1$. For $t=\Psi\binom f0$ the first block of
$\Phi t=c_p\binom f0$ is $H_tt=c_pf$.
\end{proof}

\begin{proposition}[Hodge--Tate inclusion for linear data]\label{pf:sh-local:lindatHT}
Let $H$ be a Hodge--Tate datum for a Hodge--Tate setting $(O,\Z_p\otimes\OO_E,\sigma,c_p)$ and
$(A,B,\rho_G,N,(s_k))$ linear data for $H$. Then $c\,f\in\langle s_1,\dots,s_m\rangle_O$ for all $f\in O^d$ and all
$c\in O$ with $v(c)^{p-1}\le v(p)$.
\formallabel{GlobalTheory.T.linearData_hodgeTate_inclusion}
\formallabel{GlobalTheory.ChF.linDat_hodgeTate_inclusion}
\formallabel{GlobalTheory.ChF.smul_mem_span_of_split}
\end{proposition}

\begin{proof}
$O$ is a domain, $c_p\ne0$ and the matrix $E$ of $H$ (with entries in $O$) is alternating with unit determinant
(\lref{ReciprocityTheory.PGsh.Up.htd_E}), so \cref{pf:sh-local:htsurj} applies. From $v(c)^{p-1}\le
v(p)=v(c_p)^{p-1}$ we get $v(c)\le v(c_p)$, so $c=ec_p$ with $e\in O$; it suffices to show
$c_pf\in\mathcal S:=\langle s_k\rangle$. Let $A',B'$ be as in the split exactness and, for $i\le r$, $s_{k(i)}$ with
$A^{\mathsf T}s_{k(i)}=\delta_i$. Then $u=\sum_i(A^{\mathsf T}f)_is_{k(i)}\in\mathcal S$ and
$A^{\mathsf T}(f-u)=0$. Transposing $AA'+B'B=1$ gives $f-u=B^{\mathsf T}B'^{\mathsf T}(f-u)$. By
\cref{pf:sh-local:htsurj} there is $t$ with $H_tt=c_pB'^{\mathsf T}(f-u)$, so
$c_p(f-u)=B^{\mathsf T}H_tt=\sum_it_iB^{\mathsf T}(H_t)_{\cdot,i}\in\mathcal S$, as each $B^{\mathsf T}(H_t)_{\cdot,i}$ is
some $s_k$. Hence $c_pf=c_p(f-u)+c_pu\in\mathcal S$.
\end{proof}

\begin{proof}[Proof of \cref{pf:sh-local:Hp}]
Given the data, \cref{pf:sh-local:lindatex} (with $\OO_E=\Z$) gives characters $u_1,\dots,u_m$ of $F^\iota$ and
linear data with $s_k=du_k$; \cref{pf:sh-local:lindatHT} gives the inclusion.
\end{proof}

\subsubsection{The unramified bound}

\begin{proposition}[Unramified length bound for linear data]\label{pf:sh-local:lindatU}
Let $p\ne2$, $O$ local, $\OO_E$ as in \cref{sec:pf:sh-local:chartsetting} with $p$ not dividing its discriminant,
$H$ a Hodge--Tate datum for $(O,\Z_p\otimes\OO_E,\sigma,c_p)$ and $(A,B,\rho_G,N,(s_k))$ linear data for $H$. Let
$w,n,t$ be as in $(\mathrm U_p)$ (Frobenius ordering $t_{j+1}(a)\equiv t_j(a)^p$), $r:\Hom(\OO_E,O)\to\N$ and
$W(\tau)$ eigenbases of $\rho_G$ of rank $r(\tau)$ at every $\tau$. Then for every $j$ there is $x$ such that
$O^{r(t_j)}/\langle\mathrm{res}_{W(t_j)}(s_k)\rangle$ has valuation length $x$ and
$v(p)^{\sum_{l\ge0}p^{-(l+1)}r(t_{j+l+1}\circ*)}\le x$.
\formallabel{GlobalTheory.T.linearData_unramified_length_bound}
\formallabel{GlobalTheory.ChF.linDat_unram_bound}
\end{proposition}

\begin{proof}
For an $O$-linear representation $\rho'$ of $\OO_E$ and $\tau:\OO_E\to C$ let $e(\rho',\tau)$ be the dimension over
$C$ of the common $\tau$-eigenspace of the $\rho'(a)$. The steps are as follows.
\begin{enumerate}[label=(\arabic*)]
\item \emph{Unramified data} (\lref{GlobalTheory.ChF.exists_unramData}): since $p\nmid\mathrm{disc}\,\OO_E$ and
$P=P_L=0$, there are a finite index set $I$ in bijection $i\mapsto\tau(i)$ with $\Hom(\OO_E,O)$, orthogonal idempotent
matrices $\eps_T(i)\in M_{2g}(O)$ and $\eps_L(i)\in M_g(O)$ with $\sum_i\eps_T(i)=1=\sum_i\eps_L(i)$,
$\rho_H(1\otimes a)\eps_T(i)=\tau(i)(a)\eps_T(i)$ and $\rho_L(1\otimes a)\eps_L(i)=\tau(i)(a)\eps_L(i)$, an involution
$i\mapsto\bar\imath$ with $\tau(\bar\imath)=\tau(i)\circ*$, a permutation $\phi$ with
$v(\tau(\phi i)(a)-\tau(i)(a)^p)<1$ for all $a$, distinct indices having embeddings that differ modulo the maximal
ideal, and ranks $r_U(i)=\rank\eps_L(i)=e(\rho_L,\tau(i))$. For $i,j\in I$
let $\mathrm{cok}(i,j)$ be the image of $\eps_L(i)^{\mathsf T}$ modulo its intersection with $H_t$ applied to the image
of $\eps_T(j)$.
\item \emph{Lengths} (\lref{GlobalTheory.ChF.hasLen_exists_mono_fp}): every finitely presented $O$-module has a
valuation length, and if $M\to M'$ is surjective then the length of $M$ is at most that of $M'$.
\item \emph{Ranks.} The local factor $A_pw$ is monogenic, $A_pw=\Z_p[\alpha]w$
(\lref{GlobalTheory.PlaceMono.place_monogenic}), and the image of $\rho_H(w)$ on $\Z_p^{2g}$ is free over it, of
some rank $h$ (\lref{GlobalTheory.LocalFactorFree.exists_free_generators}); hence every embedding $\tau$ belonging
to $w$ has $e(\rho_H,\tau)=h$ (\lref{GlobalTheory.ChF.eigDim_eq_of_basis}), and
$e(\rho_H,\tau)=e(\rho_L,\tau)+e(\rho_L,\tau\circ*)$ (the Hodge--Tate decomposition of the Tate module,
\lref{GlobalTheory.ChF.eigDim_tate}). So $r_U(i)+r_U(\bar\imath)=h$ for the indices $i$ of embeddings belonging
to $w$. By split exactness, eigenspace dimensions are additive:
$e(\rho_G,\tau)=e(N,\tau)+e(\rho_L,\tau)$ (\lref{GlobalTheory.EigAdd.eigDim_add_transpose}), and
$e(\rho_G,\tau)=r(\tau)$ for an eigenbasis of rank $r(\tau)$
(\lref{GlobalTheory.ChF.eigB_res_surjective}); hence $r_U(i)\le r(\tau(i))\le d$.
\item \emph{Comparison of cokernels.} Let $i_j$ be the index of $t_j$. The module $\mathrm{cok}(i_j,\bar\imath_j)$ has a
valuation length $y_1$ (\lref{ReciprocityTheory.PGsh.Up.lenEx}), and there is a surjection from it onto
$O^{r(t_j)}/\langle\mathrm{res}_{W(t_j)}(s_k)\rangle$ (\lref{GlobalTheory.ChF.cokE_surj_res_quot}); so
$y_1\le x$ by (2), where $x$ is the valuation length of the latter module.
\item \emph{The Frobenius bound.} The unramified cokernel length theorem
(\lref{ReciprocityTheory.PGsh.Up.u7}) states: for $p\ne2$, if $\mathrm{cok}(i,\bar\imath)$ and $\mathrm{cok}(\bar\imath,i)$
have valuation lengths $V_1$ and $V_2$, then
\[
V_1=V_2\cdot v(p)^{\frac{r_U(i)}{p-1}-\sum_{l\ge0}p^{-(l+1)}r_U(\phi^{l+1}i)} .
\]
The module $\mathrm{cok}(\bar\imath,i)$ is generated by $r_U(\bar\imath)$ elements
(\lref{ReciprocityTheory.PGsh.Up.idemGen}) and killed by $c_p$ (\lref{ReciprocityTheory.PGsh.Up.up_kill}),
so $V_2\ge v(c_p)^{r_U(\bar\imath)}=v(p)^{r_U(\bar\imath)/(p-1)}$
(\lref{ReciprocityTheory.PGsh.Up.killedGen}). Hence
$v(p)^{\frac{r_U(i)+r_U(\bar\imath)}{p-1}-\sum_{l\ge0}p^{-(l+1)}r_U(\phi^{l+1}i)}\le y_1$ for $i=i_j$
(\lref{ReciprocityTheory.PGsh.Up.up_full}).
\item \emph{Exponents.} With $h=r_U(\phi^{l+1}i_j)+r_U(\overline{\phi^{l+1}i_j})$ for all $l$ (by (3), as all
$\phi^{l+1}i_j$ belong to $w$) and $r_U(\overline{\phi^{l+1}i_j})\le r(t_{j+l+1}\circ*)\le d$, the elementary
inequality \lref{GlobalTheory.ChF.exp_bound}
\[
\frac h{p-1}-\sum_{l\ge0}p^{-(l+1)}a_l\le\sum_{l\ge0}p^{-(l+1)}c_l\qquad(a_l+b_l=h,\ b_l\le c_l\le d)
\]
(a consequence of $\sum_{l\ge0}p^{-(l+1)}=\frac1{p-1}$) shows that the exponent in (5) is at most
$\sum_lp^{-(l+1)}r(t_{j+l+1}\circ*)$; since $v(p)<1$, the claim follows from (4) and (5).
\end{enumerate}
\end{proof}

\begin{proof}[Proof of \cref{pf:sh-local:Up}]
By \cref{pf:sh-local:lindatex} there are characters $u_k$ of $F^\iota$ and linear data with $s_k=du_k$ and
$\rho_G=\rho$; an eigenbasis for the action on $F^\iota$ is an eigenbasis of $\rho$
(\lref{GlobalTheory.ChF.isEigBasis_iff}). \Cref{pf:sh-local:lindatU} gives $x$, and the characters $u_k$ witness
$\mathrm{AF}(W(t_j),x)$.
\end{proof}

\subsubsection{The split-place bound}

\begin{proposition}[Split-place length bound for linear data]\label{pf:sh-local:lindatA}
Let $H$ be a Hodge--Tate datum for $(O,\Z_p\otimes\OO_E,\sigma,c_p)$, $(A,B,\rho_G,N,(s_k))$ linear data for $H$, and
$w,n,t,v_d,r_p,r_m,(W_j)$ as in $(\mathrm A_p)$, with eigenbases of $\rho_G$. Then there are $x_1,\dots,x_n$ such that
$O^{r_p}/\langle\mathrm{res}_{W_j}(s_k)\rangle$ has valuation length $x_j$ and
$v(p)^{n\,r_m(\frac1{p-1}+\frac{v_d}2)}\le\prod_jx_j$.
\formallabel{GlobalTheory.T.colength_bound_split_place_linear_data}
\formallabel{GlobalTheory.ChF.colength_bound_linDat}
\end{proposition}

\begin{proof}
\begin{enumerate}[label=(\arabic*)]
\item $A_pw=\Z_p[\alpha]w$ has the $\Z_p$-basis $(w\alpha^k)_{k<n}$
(\lref{GlobalTheory.PlaceMono.place_monogenic}); the values $a_j=\tilde t_j(\alpha)$ are pairwise distinct and
their discriminant $\Delta_0=\prod_{j<j'}(a_j-a_{j'})^2$ has valuation $\Delta:=v(\Delta_0)=v(p)^{nv_d}$, with
$nv_d\ge0$ (\lref{GlobalTheory.PowerBasisDiscr.discr_power_basis}).
\item For $\tau\in\mathrm{Emb}(\bar w)$, $e(\rho_L,\tau)\le r_m$ (additivity of eigenspace dimensions as in the proof
of \cref{pf:sh-local:lindatU}).
\item \emph{The abelian part} (\lref{GlobalTheory.T.abelianPart_cokernel_length_le}): the cokernel module of the
Hodge--Tate datum between $w$ and $\bar w$ (formed with $H_t$) has a valuation length $c_w$ with
$v(p)^{nr_m/(p-1)}\le c_w$.
\item \emph{Generators} (\lref{GlobalTheory.T.exists_abelian_torus_generators_at_place}): there are $h\le r_p+r_m$
and vectors $\psi_1,\dots,\psi_h\in O^d$ such that each $\mathrm{res}_{W_j}(\psi_l)$ lies in
$\langle\mathrm{res}_{W_j}(s_k)\rangle$, and a surjection from the module of (3) onto $M/M'$, where
$M\subseteq\prod_jO^{r_p}$ is the image of $\kappa:f\mapsto(\mathrm{res}_{W_j}(f))_j$ and $M'$ is spanned by the vectors
$(\tilde t_j(\alpha)^k\,\mathrm{res}_{W_j}(\psi_l))_j$, $k<n$, $l\le h$. $M$ is stable under multiplication by
$(\tilde t_j(\alpha))_j$ and surjects onto each factor (\lref{GlobalTheory.ChF.eigB_res_surjective}). By (2) of the
previous proof, $M/M'$ has a valuation length $c$ with $c_w\le c$.
\item \emph{The lattice inequality} (\lref{GlobalTheory.T.lattice_inequality_valuation_integers}): for such
$M,M'$ there are $y_1,\dots,y_n$ such that $O^{r_p}/\langle\mathrm{res}_{W_j}(\psi_l)\rangle_l$ has valuation length
$y_j$ for every $j$, and $c^2\,\Delta^h\le(\prod_jy_j)^2\,\Delta^{r_p}$, with $\Delta$ as in (1).
\item Since $\langle\mathrm{res}_{W_j}(\psi_l)\rangle\subseteq\langle\mathrm{res}_{W_j}(s_k)\rangle$, the valuation
length $x_j$ of $O^{r_p}/\langle\mathrm{res}_{W_j}(s_k)\rangle$ satisfies $y_j\le x_j$. The numerical lemma
\lref{GlobalTheory.ChF.lav_num} combines $v(p)^{nr_m/(p-1)}\le c_w\le c$, $c^2\Delta^h\le(\prod y_j)^2\Delta^{r_p}$,
$h\le r_p+r_m$ and $\Delta=v(p)^{nv_d}\le1$ into $v(p)^{nr_m(\frac1{p-1}+\frac{v_d}2)}\le\prod_jy_j\le\prod_jx_j$.
\end{enumerate}
\end{proof}

\begin{proof}[Proof of \cref{pf:sh-local:Ap}]
As for \cref{pf:sh-local:Up}, with \cref{pf:sh-local:lindatA} in place of \cref{pf:sh-local:lindatU}.
\end{proof}

% ---------------------------------------------------------------------------
\subsection{The period tower theorem}\label{sec:pf:sh-local:ptd}

We use the notation of \cref{sec:pf:tower}: $p$ is a prime, $C$ an algebraically closed field of characteristic $0$
with a valuation $v$ and ring of integers $O$ ($p$-adically complete, $p\notin O^\times$); $\eps\in O^\flat$ is such
that $(\eps^{1/p})^\sharp$ is a primitive $p$-th root of unity, $A_{\inf}=W(O^\flat)$, $\xi\in A_{\inf}$ the
corresponding generator of the kernel of Fontaine's map $\theta:A_{\inf}\to O$, and $R_0=A_{\inf}/\xi^2$, a local ring
(\lref{GlobalTheory.WdAinf.ainf_thetaQuotRing_isLocalRing}). The \emph{$A_{\inf}$-tower} is the tower of rings
$A_{\inf}/(\xi^2,p^{n+1})\to O/p^{n+1}$ together with its limit ring $O_\infty$ (an $R_0$-algebra which is again the
ring of integers of $v$). For a local $R_0$-algebra $R$ we write
$G(R)$ for the group of $R$-points of a group scheme $G$ over $R_0$.

\begin{theorem}[Period tower datum]\label{pf:sh-local:ptd}
Let $g\in\N$ and let $G$ be a commutative group scheme over $R_0$ and $P$ a line bundle on $G\times G$ such that: $G$
is an abelian scheme, smooth of relative dimension $g$; $P$ has a generating section of its cube bundle $\Theta(P)$ and
is trivial on $0\times G$ and on $G\times0$; $P$ is isomorphic to its pull-back by the swap of the factors; the
multiplication by $p$ on $G$ is finite, flat and surjective; every $C$-point $x$ of $G$ for which the pull-back of
$P$ along $y\mapsto(y,x)$ is trivial is $0$; and $G(C)$ has exactly $p^{2g}$ points of order dividing $p$. Let $G'$ be
an abelian scheme over $O_\infty$ with a morphism $\varphi:G'\to G$ exhibiting $G'$ as the base change of $G$ (as a
group), fix a $\Z_p$-algebra structure on $O_\infty$, and let $O_0$ be a commutative ring with an involution $a\mapsto
a^*$ acting on $G'$ by endomorphisms $\rho(a)$ such that $\rho(a^*)$ is the Rosati adjoint of $\rho(a)$ with respect to
$P$ (the line bundle $P(\rho(a)\times1)\otimes P(1\times\rho(a^*))^{-1}$ on $G'\times G'$ is trivial). Then there are:
\begin{enumerate}[label=(\roman*)]
\item a functor $B$ on local $R_0$-algebras with values in abelian groups, endowed with the structure of a biextension
functor, and natural additive bijections $e_R:B(R)\to G(R)$;
\item for each $a\in O_0$ an adjoint pair $\mathrm{Ad}(a)$ of endomorphisms of $B$ over $O_\infty$ inducing $\rho(a)$
through $e$ and $\varphi$;
\item a dual tower of the $A_{\inf}$-tower and a period tower datum $\mathcal D$ (in the sense described in
\cref{sec:pf:tower}, for the ring $\Z_p\otimes O_0$ with the involution $1\otimes*$ and the constant $c_p$ of the
tower) whose Hodge--Tate datum has $P=P_L=0$, such that $\mathrm{Ad}(a)$ acts on tangent vectors at the origin through
$\rho_L(1\otimes a)$ in the formal chart basis of $\mathcal D$, and on the torsion points
$\tau_j(n)$ of $\mathcal D$ by $\mathrm{Ad}(a)\tau_j(n)=\sum_k\overline{\rho(1\otimes a)_{kj}}\,\tau_k(n)$, the bar
denoting reduction modulo $p^{n+1}$.
\end{enumerate}
\formallabel{GlobalTheory.T.exists_periodTowerDatum_of_abelianScheme_symmetricBundle}
\formallabel{GlobalTheory.AbelianPeriodDatum.main}
\end{theorem}

\begin{proof}
The formal proof assembles the following steps, all of which are results of the development.
\begin{enumerate}[label=(\arabic*)]
\item \emph{Biextension.} The theorem of the cube for $P$ (a generating section of $\Theta(P)$) and the trivialisations
along $0\times G$, $G\times0$ give $P$ the structure of a biextension of $G\times G$ by $\mathbb G_m$
(\lref{GlobalTheory.nonempty_isBiext_of_cube_gen}); by the symmetry of $P$ it is symmetric, with a generating
section of $P\otimes\mathrm{swap}^*P^{-1}$ compatible with the two partial group laws
(\lref{GlobalTheory.BiextSym.main}).
\item \emph{Points.} Since $R_0$ is local and $G$ is smooth of relative dimension $g$, the points of $G$ with values in
local $R_0$-algebras are described by nilpotent coordinates and the fibre-product property of points
(\lref{GlobalTheory.NilCoord.nilpotent_coordinates}, \lref{GlobalTheory.points_fibre_product}); together with
the factor systems of the biextension this gives the functor $B$ with the bijections $e_R$
(\lref{GlobalTheory.BXOf.main}, \lref{GlobalTheory.GroupPointsBX.main},
\lref{GlobalTheory.exists_BX_of_factorSystems}).
\item \emph{Adjoint pairs.} The Rosati adjointness of $\rho(a)$ and $\rho(a^*)$ gives, for each $a$, an adjoint pair of
endomorphisms of $B$ over $O_\infty$ for the biextension pairing (\lref{GlobalTheory.exists_oAdj_of_rosatiAdjoint}),
and these satisfy the ring laws $\mathrm{Ad}(1)=1$, $\mathrm{Ad}(a+b)=\mathrm{Ad}(a)+\mathrm{Ad}(b)$,
$\mathrm{Ad}(ab)=\mathrm{Ad}(a)\mathrm{Ad}(b)$ and $\mathrm{Ad}(a)^{\mathrm{adj}}=\mathrm{Ad}(a^*)$
(\lref{GlobalTheory.AbelianPeriodDatum.ring_laws}).
\item \emph{Torsion.} $B(O_\infty)$ is $p$-divisible, and its $p$-torsion has as many elements as $G(C)$, i.e.
$p^{2g}$ (\lref{GlobalTheory.AbelianPeriodDatum.divisible_and_card}, from the finiteness and surjectivity of
multiplication by $p$ and the description of points over valuation rings
\lref{GlobalTheory.points_divisible_and_torsion_of_valuation_integers}).
\item \emph{Non-degeneracy.} The Weil pairing of the biextension on the $p$-torsion of $B(O_\infty)$ is non-degenerate
(\lref{GlobalTheory.pairing_nondegenerate_integral_torsion}), using the non-degeneracy on the geometric fibre
(\lref{GlobalTheory.biext_pairing_nondegenerate_geometric_fibre}), which rests on the descent of invariant
generators along multiplication by $N$ (\lref{GlobalTheory.invariantGenerator_descends_mulN}) and the
generation of torsion (\lref{GlobalTheory.TorsionGen.main}).
\item \emph{The period tower datum} (\lref{GlobalTheory.PeriodDatum.main}). From (2)--(5): Witt vector data for
the tower (\lref{GlobalTheory.wData_ainfTower_exists}, with the Teichm\"uller quotient
\lref{GlobalTheory.WittQuot.teichmuller_quot}), a formal chart of $B$ at the origin
(\lref{GlobalTheory.FormalChartOf.exists_formalChart}), the lifting of points along nilpotent thickenings
(\lref{GlobalTheory.SmoothNilLift.main}), a finite flat model of the $p$-torsion
(\lref{GlobalTheory.FinOf.main}), a dual tower (\lref{GlobalTheory.exists_dualTower}), a compatible system of
bases $\tau_j(n)$ of the $p^{n+1}$-torsion (\lref{GlobalTheory.PeriodDatum.exists_compatible_basis}), the
differentials of the adjoint pairs as a ring homomorphism (\lref{GlobalTheory.exists_endDifferential_ringHom}),
and the Tate matrices linking heights, the Weil pairing and the action
(\lref{GlobalTheory.exists_tateMatrices_link}, with \lref{GlobalTheory.tateBasis_gram} and
\lref{GlobalTheory.PadicOrderAction.main}) are assembled into the datum $\mathcal D$ with $P=P_L=0$ and the two
compatibilities of (iii).
\end{enumerate}
\end{proof}

% ---------------------------------------------------------------------------
\subsection{Moduli schemes and the theta-constant formula}\label{sec:pf:sh-local:moduli}

\subsubsection{Moduli schemes}
\begin{definition}[Data, levels, objects and moduli schemes]\label{pf:sh-local:def:moduli}
Let $E$ be a CM field with a toroidal host $H$ (\cref{pf:sh-local:def:torhost}). A \emph{datum} $D_w$ for $H$ is a
finitely generated projective $\OO_E$-module $\Lambda$ of rank $p+q$ with a non-degenerate alternating $\Z$-bilinear
form $\langle\,,\rangle$ satisfying $\langle ax,y\rangle=\langle x,\bar ay\rangle$, together with a complex structure
$J$ on $\Lambda\otimes\R$ preserving the form, with $\langle Jx,x\rangle>0$ for $x\ne0$, commuting with $\OO_E$ and
satisfying the signature condition for $(p,q)$ relative to $\Psi_H$; its \emph{dimension} is
$\dim D_w=\frac12[E:\Q]\,pq$. A \emph{level} $l$ consists of an odd integer $M\ge3$, an ideal $\mathfrak M\subseteq\OO_E$
generated by elements of $\Z[w_0]$ with $\mathfrak M\bar{\mathfrak M}=M\OO_E$ and $\mathfrak M+\bar{\mathfrak M}=\OO_E$,
and an integer $l_A\ge1$. Let $K_l=\Q(\mathrm{rt}_0,\zeta_{l_A})\subset\C$. An \emph{object} of type $(D_w,l)$ over a
$K_l$-scheme $S$ is an abelian scheme $G$ over $S$ (a proper semi-abelian scheme) with an action of $\OO_E$, a
rigidified symmetric line bundle with cubical structure which is relatively ample and has trivial kernel (a principal
polarisation), for which the action is Rosati-compatible, together with level structures $\eta:\Lambda\to
G[\mathfrak M]$ and $\eta_A:\Lambda\to G[l_A]$, such that at every geometric point the signature condition holds and
there is a Tate isomorphism compatible with the Weil pairings, with $\zeta_{l_A}$ and with $\eta,\eta_A$. A
\emph{moduli scheme} of type $(D_w,l)$ is a scheme $X$ quasi-projective over $K_l$ and smooth of relative dimension
$\dim D_w$, with an object $\mathcal U$ over $X$, such that for every \emph{reduced} $K_l$-scheme $S$ every object over
$S$ is the pull-back of $\mathcal U$ along a unique morphism $S\to X$, and such that a dual-number point of $X$ whose
object is isomorphic to the pull-back of $\mathcal U$ along the underlying constant point is constant.
\formal{ReciprocityTheory.ChartRestriction.G.Geo.Tor.Datum}
\formal{ReciprocityTheory.ChartRestriction.G.Geo.Tor.Level}
\formal{ReciprocityTheory.ChartRestriction.G.Geo.Tor.Obj}
\formal{ReciprocityTheory.ChartRestriction.G.Geo.Tor.Moduli}
\end{definition}

\begin{theorem}[Moduli schemes exist]\label{pf:sh-local:moduli}
For every CM field $E$ with toroidal host $H$, every datum $D_w$ for $H$ with $p+q\ge1$ and every level $l$ there is a
moduli scheme of type $(D_w,l)$.
\formal{ReciprocityTheory.ChartRestriction.G.Geo.Tor.N1ab}
\formallabel{GlobalTheory.Moduli.moduli_scheme_exists}
\formallabel{GlobalTheory.T.moduliScheme_nonempty_at_every_level}
\formallabel{GlobalTheory.ZM.moduli_nonempty_all_levels}
\end{theorem}

\begin{remark}
This is a version of the existence of fine moduli schemes of PEL type: for principally polarised abelian varieties
with full level $n\ge3$ structure it is \cite[Theorem~7.9]{MumfordFogartyKirwan1994}, and in the PEL setting it is in
\cite[\S5]{Kottwitz1992}, \cite[Ch.~1]{Lan2013}. The formal statement is weaker in that the universal property is
required only for reduced test schemes (with a separate rigidity condition on dual-number points), and it is
stated over the field $K_l$ only; it is proved in the development from the laws of abelian schemes and the theta
uniformisation of complex abelian varieties, not from geometric invariant theory.
\end{remark}

\begin{proof}
The proof uses the laws of abelian schemes
(\lref{GlobalTheory.T.all_six_abelian_scheme_laws_hold}), the theorem of the cube for polarised abelian schemes in
the form ``every line bundle has a generating section of its cube bundle''
(\lref{GlobalTheory.T.theorem_of_cube_polAbelianScheme}) and the theta uniformisation of complex abelian varieties
(\lref{GlobalTheory.T.theta_uniformisation_complex_abelian_varieties}). Let $E,H,D_w,l$ be given.
\begin{enumerate}[label=(\arabic*)]
\item There is an odd prime $n$ such that for every $A_0\ge1$ there are levels $l_0,l_2$ with $M(l_0)=n$, $l_A(l_0)=A_0$,
$l_0\mid l_2$ and $l\mid l_2$ (\lref{GlobalTheory.T.exists_level_with_odd_prime_first_number}).
\item \emph{Prime-power level.} There is $j$ such that a moduli scheme exists for every level $l_0$ with $M(l_0)=n$ and
$l_A(l_0)=(n^j)^2$ (\lref{GlobalTheory.T.exists_moduli_prime_power_fullLevel}; it is constructed from a family of
abelian varieties with level structure over a finite type base, \lref{GlobalTheory.ZM.FullLevel.moduli_of_family}).
\item \emph{Passing to a multiple.} If $l_0\mid l_2$ and a moduli scheme exists at level $l_0$, one exists at level $l_2$
(\lref{GlobalTheory.T.nonempty_moduli_of_level_dvd}, via a finite \'etale cover adding level structure,
\lref{GlobalTheory.T.exists_finiteEtale_deeper_level_cover}).
\item \emph{Passing to a divisor.} If $l\mid l_2$ and a moduli scheme exists at level $l_2$, one exists at level $l$
(\lref{GlobalTheory.T.nonempty_cmModuli_of_deeper_level}: the quotient of the moduli scheme of level $l_2$ by the
finite group of level changes exists, is unramified on dual-number points and represents the moduli problem on reduced
schemes, \lref{GlobalTheory.T.nonempty_levelQuot_of_level_dvd},
\lref{GlobalTheory.T.levelQuotient_unramified_dualNumberPoints},
\lref{GlobalTheory.T.levelQuotient_represents_reduced}).
\end{enumerate}
Apply (1) with $A_0=(n^j)^2$, then (2) at $l_0$, (3) to get level $l_2$, and (4) to get level $l$.
\end{proof}

\subsubsection{The theta-constant formula}
\begin{definition}[Key inputs]\label{pf:sh-local:def:keyin}
A \emph{key input} $x$ over a commutative ring $R$ consists of a polarised abelian scheme $P$ over $R$ of relative
dimension $g\ge1$ with line bundle $L$, a generating section of $L$ along the unit section (a rigidification), a section
$\theta$ of $L$ which is non-zero at some point of every fibre, and a basis $v_1,\dots,v_g$ of the Lie algebra (tangent
vectors at the origin, as morphisms from $\Spec R[\epsilon]$). For a ring homomorphism $f:R\to R'$, a key input $x'$
over $R'$ is a \emph{pull-back} of $x$ along $f$ if there is a morphism $P'\to P$ over $\Spec f$ exhibiting $P'$ as the
base change of $P$ as a group scheme, together with a generating isomorphism of $L'$ with the pull-back of $L$, which
carries the rigidification of $x'$ to that of $x$, $\theta'$ to the pull-back of $\theta$, and $v'_i$ to $v_i$.
\formal{ReciprocityTheory.ChartRestriction.G.Geo.Tor.KeyIn}
\formal{ReciprocityTheory.ChartRestriction.G.Geo.Tor.KeyPB}
\end{definition}

Over $\C$, a \emph{holomorphic uniformisation}
of $P$ is a surjective homomorphism $\C^g\to P(\C)$ with kernel a lattice $\Lambda$; for such a uniformisation, a
Hermitian form $H_m$ and a semi-character $\chi$ defining $L$ (Appell--Humbert data), the section $\theta$ corresponds
to a theta function $\vartheta$ on $\C^g$ and the $v_i$ to vectors $W_i\in\C^g$. Let
\[
\vartheta_{a,b}(\tau,z)=\sum_{n\in\Z^g}\exp\bigl(\pi i(n+a)^{\mathsf T}\tau(n+a)+2\pi i(n+a)^{\mathsf T}(z+b)\bigr)
\]
be the theta function with characteristic $(a,b)$.

\begin{definition}[The theta-constant formula]\label{pf:sh-local:def:key}
$\sfname{KeyFormula}$ asserts that there is a function $\kappa$ assigning to every field $K$ of characteristic $0$ and
every key input $x$ over $K$ a unit $\kappa(x)\in K^\times$ such that:
\begin{enumerate}[label=(\roman*)]
\item \emph{(transformation law)} if $x'$ differs from $x$ only by $\theta'=a\theta$ ($a\in K$) and Lie basis
$v'_i=\sum_jM_{ij}v_j$, then $a^8\kappa(x')=\det(M)^4\kappa(x)$;
\item \emph{(base change)} $\kappa(x')=f(\kappa(x))$ whenever $x'$ is the pull-back of $x$ along $f:K\to K'$;
\item \emph{(integrality)} if $x$ is a key input over a discrete valuation ring $R$ with fraction field $K$ and $x_K$ its
pull-back to $K$, then $\kappa(x_K)$ is a unit $u$ of $R$, and $\kappa(x_k)=r(u)$ for every surjection $r$ of $R$ onto a
field $k$ of characteristic $0$ and every pull-back $x_k$ of $x$ along $r$;
\item \emph{(normalisation)} if $x$ is a key input over $\C$ with a holomorphic uniformisation with lattice
$\Z^g+\tau\Z^g$ ($\tau$ symmetric with positive definite imaginary part), compatible with the rigidification, under
which $\theta$ is $\vartheta_{a,b}(\tau,\cdot)$ with $a,b\in\{0,\frac12\}^g$ and $v_i=(2\pi i)^{-1}\partial/\partial z_i$,
then $\kappa(x)=1$;
\item \emph{(norm formula)} for every key input $x$ over $\C$ with a holomorphic uniformisation with Appell--Humbert data
$(H_m,\chi)$, theta function $\vartheta$ and vectors $W_i$ as above, and every fundamental domain $\mathcal F_\Lambda$
of $\Lambda$,
\[
|\kappa(x)|\,n_\vartheta^4\,n_\eta^2=(2\pi)^{2g},\qquad
n_\vartheta=\frac1{\vol\mathcal F_\Lambda}\int_{\mathcal F_\Lambda}|\vartheta(z)|^2e^{-\pi H_m(z,z)}\,dz,\quad
n_\eta=\frac{\vol\mathcal F_\Lambda}{\pi^g|\det W|^2}.
\]
\end{enumerate}
\formallabel{ReciprocityTheory.ChartRestriction.G.Geo.Tor.KeyFormula}
\end{definition}

\begin{theorem}[The theta-constant formula]\label{pf:sh-local:key}
$\sfname{KeyFormula}$ holds.
\formallabel{GlobalTheory.KeyFam.keyFormula_holds}
\end{theorem}

\begin{remark}
In terms of the Hodge bundle $\omega$ and the determinant of $\pi_*L$, property (i) says that $\kappa$ is a section of
$(\pi_*L)^{\otimes-8}\otimes\omega^{\otimes4}$ (up to the conventions for $\det$), and (ii)--(v) say that it is a
trivialisation compatible with base change, integral, and with the stated archimedean norm. This is Moret-Bailly's
\emph{formule cl\'e} \cite{MoretBailly1985} for principally polarised abelian varieties, in the form used for heights;
the formal statement is restricted to bases which are fields of characteristic $0$ and discrete valuation rings (with
residue fields of characteristic $0$ in (iii)), and the formal proof is independent of \cite{MoretBailly1985}.
\end{remark}

\begin{proof}
The formal proof applies \lref{GlobalTheory.keyFormula_of_connFam} to
\lref{GlobalTheory.KeyFam.connFam_of_polFam}, applied in turn to \lref{GlobalTheory.KeyFam.polFam}, with the
laws of abelian schemes (\lref{GlobalTheory.abelianSchemeLaws_all}), the theta uniformisation
(\lref{GlobalTheory.thetaUnifAV_closed}) and the algebraicity of theta tori
(\lref{GlobalTheory.torusAlgebraic_holds}: for every $\tau$ and characteristic $(a,b)\in\{0,\frac12\}^{2g}$ there is a
complex principally polarised abelian variety uniformised by $\Z^g+\tau\Z^g$ on which $\vartheta_{a,b}(\tau,\cdot)$ is a
section, with very ample powers and all endomorphisms of the lattice realised). It has three steps.

\emph{(a) Polarised families} (\lref{GlobalTheory.KeyFam.polFam}). Every polarised abelian variety $P$ over $\C$ of
dimension $g\ge1$ is the fibre at a complex point $s$ of a polarised abelian scheme $Q$ over a geometrically integral
algebra $R$ over a number field $K$ which also has a $K$-point $s_0$, such that every fibre of $Q$ at $s_0$ is
\emph{potentially known} (its theta values are known after a finite extension). The construction uses a Siegel family
over a number field with level structure (\lref{GlobalTheory.T.exists_siegelFamily_over_numberField}), a point of
it over a number field with a known theta value (\lref{GlobalTheory.T.exists_numberField_point_siegelFamily_knownValue}),
the connectedness of the family through standard level points
(\lref{GlobalTheory.T.siegelFamily_standard_level_points_same_component}), a finite \'etale cover with a
symmetric square root of the principal polarisation (\lref{GlobalTheory.T.exists_symmetric_sqrt_principal_siegel_finite_etale}),
a geometrically integral piece of an \'etale cover containing both points
(\lref{GlobalTheory.T.geom_integral_of_etale_faithfully_flat_cover}), and the transfer of known values between
abelian varieties with the same $\lambda$-bundle (\lref{GlobalTheory.T.potentiallyKnown_of_same_lambdaBundle}).

\emph{(b) Connected families of key inputs} (\lref{GlobalTheory.KeyFam.connFam_of_polFam}). From (a): every key
input $x$ over $\C$ is the pull-back along some $s:S\to\C$ of a key input $y$ over a smooth algebra $S$ over a number
field $k$ which stays a domain after every field extension, and there are a $k$-point $s_0$ of $S$ and $t_0\in k$ such
that for every embedding $e:k\to\C$, every pull-back $x_0$ of $y$ along $e\circ s_0$ has a normal form
$(\tau,M,a,b,c)$ --- a holomorphic uniformisation with lattice $\Z^g+\tau\Z^g$ under which $\theta=c\,\vartheta_{a,b}(\tau,\cdot)$,
$(a,b)\in\{0,\frac12\}^{2g}$, $c\ne0$, and $v_i=(2\pi i)^{-1}\sum_jM_{ij}\partial/\partial z_j$ --- with
$e(t_0)=\det(M)^4/c^8$. This uses key data near two primes (\lref{GlobalTheory.T.exists_keyData_near_two_primes}:
a Lie basis and a section of the principal polarisation non-vanishing near two given points) and universal frames on
the family (\lref{GlobalTheory.T.exists_thetaKeyFamily_universal_frames}).

\emph{(c) Definition of $\kappa$} (\lref{GlobalTheory.keyFormula_of_connFam}). Over $\C$ every key input has a normal
form, and one puts $t(x)=\det(M)^4/c^8$. Using (b), $t$ does not depend on the normal form and is equivariant under
field automorphisms of $\C$ (\lref{GlobalTheory.ThetaEquiv.val_equivariant}), and it specialises correctly in
families over local principal ideal domains (\lref{GlobalTheory.ThetaFrameSpec.value_specialises}): if $x,x'$ are
pull-backs of one key input over such a ring $S$ along $f,f':S\to\C$ with $f$ injective, and $t(x)=f(u)$, then
$t(x')=f'(u)$. Call $u\in K$ a \emph{value} of $t$ at a key input $x$ over a ring $K$ if $x$ is the pull-back along
some $i:S\to K$, with $S$ countable, of a key input $y$ over $S$, and $u=i(u_0)$ for some $u_0\in S$ such that
$t(x')=f(u_0)$ for every ring homomorphism $f:S\to\C$ and every pull-back $x'$ of $y$ along $f$. If $u$ is a value at
$x$ and $x'$ is a pull-back of $x$ along $f$, then $f(u)$ is a value at $x'$; a value at a complex key input $x$ is
$t(x)$. Since key inputs over local principal ideal domains descend to countable subrings
(\lref{GlobalTheory.T.keyData_descend_countable_subrings}), at every key input over a field $K$ of characteristic
$0$ there is a value, which is non-zero (\lref{GlobalTheory.KeyRoot.exists_isVal}) and unique
(\lref{GlobalTheory.KeyRoot.IsVal.unique}); this value is $\kappa(x)$. In
\lref{GlobalTheory.KeyRoot.keyFormula_of_families}: (i) follows from the behaviour of normal forms under scaling of
$\theta$ and change of Lie basis (\lref{GlobalTheory.KeyRoot.normalAt_scale}); (ii) from the compatibility of
values with pull-backs and uniqueness; (iv) from the normal form with $M=1$, $c=1$; (v) from the normal form attached to
the given uniformisation and the computation of its archimedean constants; and (iii) from a \emph{key rule}
(\lref{GlobalTheory.KeyTop.exists_keyRule_charHalf}): an invariant $\kappa_N$ of key inputs over fields of
characteristic $0$ which satisfies (ii), (iii), (iv) and (i) with the exponents $8N(g),4N(g)$, for some $N(g)\ge1$, and
equals $t^{N(g)}$ over $\C$. Indeed, since $\kappa_N$ is compatible with base change and equals $t^{N(g)}$ over $\C$,
$\kappa(x)^{N(g)}=\kappa_N(x)$ for every key input $x$ over a field of characteristic $0$
(\lref{GlobalTheory.KeyRoot.IsVal.pow}); by property (iii) of $\kappa_N$, $\kappa(x_K)^{N(g)}$
is the image of a unit $u'$ of the discrete valuation ring $R$. Hence $\kappa(x_K)$ is integral over $R$, so it is the
image of some $u\in R$ ($R$ is integrally closed), and $u^{N(g)}=u'$ shows that $u$ is a unit. Moreover $u$ is a value of
$t$ at $x$ itself (\lref{GlobalTheory.KeyRoot.IsVal.special}, by descending $x$ to a countable subring), so for every
surjection $r$ of $R$ onto a field $k$ of characteristic $0$, $r(u)$ is a value at $x_k$, and $\kappa(x_k)=r(u)$ by
uniqueness.
\end{proof}
